% ============================================================
%  PREAMBLE — compile with pdfLaTeX
% ============================================================
\documentclass[12pt,a4paper]{report}

% ── Encoding & Language ──────────────────────────────────────
\usepackage[utf8]{inputenc}
\usepackage[T1]{fontenc}
\usepackage[english]{babel}

% ── Page Layout ──────────────────────────────────────────────
\usepackage{geometry}
\geometry{
  a4paper,
  top=2.5cm,
  bottom=2.5cm,
  left=3cm,
  right=2.5cm
}

% ── Typography & Spacing ─────────────────────────────────────
\usepackage{setspace}
\onehalfspacing
\usepackage{parskip}
\usepackage{ragged2e}

% ── Mathematics ──────────────────────────────────────────────
\usepackage{amsmath}
\usepackage{amssymb}

% ── Graphics & Colors ────────────────────────────────────────
\usepackage{graphicx}
\usepackage{xcolor}
\usepackage{tikz}

% ── Tables ───────────────────────────────────────────────────
\usepackage{array}
\usepackage{booktabs}
\usepackage{tabularx}
\usepackage{adjustbox}

% ── Boxes & Frames ───────────────────────────────────────────
\usepackage{mdframed}

% ── Headers & Footers ────────────────────────────────────────
\usepackage{fancyhdr}

% ── Section & Chapter Formatting ─────────────────────────────
\usepackage{titlesec}
\titleformat{\chapter}[block]
  {\normalfont\Large\bfseries}{}{0em}{}
\titlespacing*{\chapter}{0pt}{0pt}{20pt}

% ── Bibliography ─────────────────────────────────────────────
\usepackage[round,authoryear]{natbib}
\newcommand{\textcite}[1]{\citet{#1}}
\newcommand{\parencite}[1]{\citep{#1}}
\bibliographystyle{plainnat}

% ── Hyperlinks ───────────────────────────────────────────────
\usepackage{hyperref}
\hypersetup{
  colorlinks=true,
  linkcolor=blue,
  citecolor=blue,
  urlcolor=blue
}

% ============================================================
\begin{document}

% ── FRONT MATTER (roman page numbers) ───────────────────────
\pagestyle{empty}
\pagenumbering{roman}

% ============================================================
%  PAGE 1 — COVER PAGE
% ============================================================
\begin{titlepage}

% Page double border
\begin{tikzpicture}[remember picture, overlay]
  \draw[line width=2pt]
    ([shift={(0.55cm,0.55cm)}]current page.south west)
    rectangle
    ([shift={(-0.55cm,-0.55cm)}]current page.north east);
  \draw[line width=0.5pt]
    ([shift={(0.7cm,0.7cm)}]current page.south west)
    rectangle
    ([shift={(-0.7cm,-0.7cm)}]current page.north east);
\end{tikzpicture}

% Watermark
\begin{tikzpicture}[remember picture, overlay]
  \node[opacity=0.07] at (current page.center)
    {\includegraphics[width=0.68\paperwidth]{picture/logo/logo.png}};
\end{tikzpicture}

  % ── Arabic header ───────────────────────────────────────
  \begin{center}
    {\small\bfseries People's Democratic Republic of Algeria}\\[1pt]
    {\small\bfseries Ministry of Higher Education and Scientific Research}
  \end{center}

  \vspace{3pt}
  \noindent\rule{\linewidth}{0.5pt}
  \vspace{4pt}

  % ── University row ──────────────────────────────────────
  \begin{center}
  \begin{tabular}{p{0.34\linewidth} c p{0.34\linewidth}}
    \raggedright
    {\small Universit\'e Badji Mokhtar -- Annaba\\
             Badji Mokhtar -- Annaba University} &
    \raisebox{-0.45cm}{%
      \includegraphics[height=1.7cm]{picture/logo/logo.png}} &
    \raggedleft
    {\small\bfseries Universit\'e Badji Mokhtar\\
                     -- Annaba --}
  \end{tabular}
  \end{center}

  \vspace{4pt}
  \noindent\rule{\linewidth}{0.5pt}
  \vspace{6pt}

  % ── Faculty block ───────────────────────────────────────
  \begin{flushleft}
    {\small \textbf{Faculty:} Technology}\\[2pt]
    {\small \textbf{Department:} Computer Science}\\[2pt]
    {\small \textbf{Domain:} Mathematics -- Computer Science}\\[2pt]
    {\small \textbf{Field:} Computer Science}\\[2pt]
    {\small \textbf{Speciality:} Machine Learning \& Intelligent Systems}
  \end{flushleft}

  \vspace{0.5cm}

  % ── Memoir label ────────────────────────────────────────
  \begin{center}
    {\large\bfseries Memoir}\\[4pt]
    {\normalsize Presented for the award of the \textbf{Master's Degree}}
  \end{center}

  \vspace{0.25cm}
  \begin{center}{\large\bfseries Theme}\end{center}
  \vspace{0.25cm}

  % ── Title box ───────────────────────────────────────────
  \begin{mdframed}[
    linecolor=black, linewidth=1pt,
    innertopmargin=14pt, innerbottommargin=14pt,
    innerleftmargin=18pt, innerrightmargin=18pt,
    roundcorner=4pt
  ]
    \begin{center}
      {\large\bfseries
        Length-Adaptive Hybrid GNN-Transformer for Sequential\\[5pt]
        Recommendation and Systematic Study of LLM Item\\[5pt]
        Description Strategies
      }
    \end{center}
  \end{mdframed}

  \vspace{0.55cm}

  % ── Author ──────────────────────────────────────────────
  \begin{center}
    \textbf{Presented by:} \quad Farouk Ouled Meriem
  \end{center}

  \vspace{0.35cm}

  % ── Supervisor ──────────────────────────────────────────
  \begin{center}
    \begin{tabular}{lll}
      \textbf{Supervisor:} & [Supervisor Name] \quad & [Grade] \quad Badji Mokhtar -- Annaba University
    \end{tabular}
  \end{center}

  \vspace{0.45cm}

  % ── Jury ────────────────────────────────────────────────
  \begin{center}
    \textbf{Defense Jury:}\\[8pt]
    \begin{tabular}{|p{3.8cm}|p{1.5cm}|p{4.2cm}|p{2.5cm}|}
      \hline
      {[Jury Member 1]} & MCB & Badji Mokhtar -- Annaba University & President   \\ \hline
      {[Supervisor]}    & MCA & Badji Mokhtar -- Annaba University & Supervisor  \\ \hline
      {[Jury Member 2]} & MCB & Badji Mokhtar -- Annaba University & Examiner    \\ \hline
    \end{tabular}
  \end{center}

  \vfill

  \begin{center}
    {\normalsize\bfseries University Year: 2024/2025}
  \end{center}

\end{titlepage}


% ============================================================
%  PAGE 2 — ACKNOWLEDGEMENTS
% ============================================================
\newpage
\chapter*{Acknowledgements}
\addcontentsline{toc}{chapter}{Acknowledgements}

\vspace{1cm}

First of all, I express my deepest gratitude to \textbf{Allah} for granting me
the strength and perseverance to accomplish this work.

\vspace{0.5cm}
I sincerely thank my \textbf{parents} for their unconditional love, support,
and sacrifices throughout my studies.

\vspace{0.5cm}
My heartfelt thanks go to my supervisor for their constant guidance,
availability, and valuable feedback at every stage of this work.

\vspace{0.5cm}
I also extend my appreciation to the members of the jury for their time
and careful evaluation of this memoir.

\vspace{0.5cm}
Finally, I thank all my colleagues and friends for their encouragement
and support throughout this journey.


% ============================================================
%  PAGE 3 — DEDICATIONS
% ============================================================
\newpage
\chapter*{Dedications}
\addcontentsline{toc}{chapter}{Dedications}

\vspace{2.5cm}

\begin{center}
  \itshape\large
  To my parents,\\[8pt]
  for every sacrifice they made so I could reach this day.\\[24pt]

  To my family and friends,\\[8pt]
  for their constant support and encouragement.\\[24pt]

  \upshape\normalsize\textbf{Farouk Ouled Meriem}
\end{center}


% ============================================================
%  TABLE OF CONTENTS, LIST OF FIGURES, LIST OF TABLES
% ============================================================
\newpage
\tableofcontents
\newpage
\listoffigures
\newpage
\listoftables


% ============================================================
%  GENERAL INTRODUCTION (Unnumbered Chapter)
% ============================================================
\newpage
\pagenumbering{arabic}  % start main matter with arabic page numbers

\chapter*{General Introduction}
\addcontentsline{toc}{chapter}{General Introduction}

% ── Section 1: Context and Motivation ──────────────────────
\section*{1. Context and Motivation}

In the era of digital transformation, a massive number of users are connected
through networks and online platforms, which have become an essential part of
daily life. E-commerce websites enable people to buy and sell products from
anywhere in the world, while streaming platforms provide instant access to
movies, music, and various forms of entertainment content.

Due to the rapid innovation occurring in the digital world, the amount of
content generated online has grown tremendously. This overwhelming abundance of
information often makes it difficult for users to make appropriate decisions or
find content that matches their interests. Faced with such a large volume of
available choices, users may experience information overload and confusion.

To address this challenge, recommendation systems have emerged as an effective
solution. These systems help users discover relevant products, services, or
content by analyzing their preferences and behaviors, thereby improving user
experience and supporting better decision-making.

Recommendation systems have various categories, one of the most important being
\textbf{sequential recommendation}, which has become a major research direction
in recent years. Unlike traditional collaborative filtering approaches that focus
on identifying user preferences based on similarities between groups of users,
sequential recommendation models emphasize the temporal order of user interactions
to predict the next item a user is likely to interact with, such as purchasing a
product, watching a movie, or listening to music. By analyzing the sequence of
past behaviors, these models can better capture dynamic user interests and provide
more accurate and personalized recommendations.

Recent breakthroughs in deep learning have significantly advanced the field of
sequential recommendation. Early models based on \textbf{Recurrent Neural Networks
(RNNs)}, such as GRU4Rec, were designed to process sequential data where the order
of user interactions plays a crucial role. These models enabled the capture of
short-term temporal dependencies in user behavior, marking an important step forward
in recommendation systems.

Later, the introduction of attention mechanisms and \textbf{Transformer-based
architectures} brought a major improvement. Models such as SASRec and BERT4Rec made
it possible to effectively model long-range dependencies within user interaction
sequences. This greatly improved both the accuracy and scalability of sequential
recommendation systems.

At the same time, \textbf{Graph Neural Networks (GNNs)} introduced a new paradigm
by modeling global relationships between users and items through the entire
interaction graph. This approach enables the extraction of rich collaborative signals
that go beyond what purely sequential models can capture, providing a more holistic
understanding of user preferences.

However, these two categories possess complementary characteristics. Graph Neural
Networks are highly effective at capturing global co-occurrence relationships within
the interaction graph, but they often show weaker performance in modeling the local
and temporal dynamics specific to individual users. On the other hand, sequential
Transformer-based models are proficient at encoding recent user preferences and
contextual dynamics, yet they generally overlook global collaborative signals hidden
within the interaction graph.

Existing works generally rely on fixed fusion strategies that combine these two
paradigms using static contribution weights. However, such approaches fail to
consider an important reality: users exhibit varying levels of interaction history,
and the contribution of each information source should adapt accordingly to the
density of this history. Users with limited interactions should rely more heavily
on the global collaborative signals provided by GNN-based models, whereas users
with rich interaction histories allow Transformer-based models to better exploit
complex sequential patterns and evolving preferences.

In parallel, the rise of \textbf{Large Language Models (LLMs)} such as GPT, LLaMA,
and BERT, pretrained on massive textual corpora, has demonstrated a remarkable
ability to encode rich and general semantic knowledge. These models have sparked a
technological revolution in the field of recommendation systems by enabling a deeper
understanding of user preferences, item semantics, and contextual information.

Their integration into recommendation systems has introduced new perspectives,
particularly for addressing the cold-start problem and enriching item representations
beyond simple identifiers. In particular, the textual embeddings generated by LLMs
from item descriptions provide powerful semantic signals that can complement the
representations learned by traditional recommendation models.

However, a core question remains largely unexplored in the literature: how should
an item's textual description be formulated in order to maximize the semantic quality
of the embeddings generated by an LLM in the context of sequential recommendation?
Indeed, different description strategies inherently encode distinct perspectives on
the item, which may critically impact the representational quality and discriminative
capacity of the resulting embeddings. Each formulation --- whether a simple title, a
metadata-enriched profile, a user-centric narrative, or a hybrid combination ---
captures a unique facet of the item and may substantially influence the relevance
and discriminative power of the derived representations.


% ── Section 2: Addressed Research Problems ─────────────────
\section*{2. Addressed Research Problems}

This memoir addresses two complementary research problems:

\medskip
\noindent\textbf{Research Problem 1 --- Adaptive GNN-Transformer Fusion Based on
Interaction History Length}

\medskip
Existing hybrid approaches that combine Graph Neural Networks and Transformer-based
architectures for sequential recommendation generally rely on static fusion
mechanisms, where the respective contributions of the two components are fixed for
all users. Such approaches ignore a fundamental disparity in the data: the length
and density of user interaction histories vary significantly from one user to another.

It is reasonable to assume that users with limited interaction histories should rely
more heavily on the global collaborative signals extracted by GNNs, which leverage
the entire interaction graph to compensate for the lack of sequential information.
Conversely, users with long and dense interaction histories provide sufficient
sequential patterns for Transformer-based models to effectively capture complex
temporal dynamics and evolving preferences. Therefore, the contribution of each
component should adapt dynamically according to the richness of the user's
interaction history.

\medskip
\noindent\textbf{Research Problem 2 --- Impact of Item Description Strategies on
LLM Embedding Quality in Sequential Recommendation}

\medskip
Existing approaches that integrate Large Language Models into sequential
recommendation pipelines generally rely on a single fixed strategy for describing
items, typically representing the item by its title or name alone. Such approaches
overlook a fundamental question: does the way an item is described to the LLM
influence the semantic quality of the resulting embeddings, and consequently the
accuracy of the recommendations?

Different description strategies inherently encode distinct perspectives on the item.
A title-only description captures the most basic identity, while richer formulations
--- such as metadata-enriched profiles, user-centric narratives, or hybrid
combinations --- may provide the LLM with more discriminative and contextually
relevant information. Therefore, the choice of description strategy may critically
impact the quality and discriminative capacity of the generated embeddings.

Despite its importance, this phenomenon has never been systematically studied.
Although Boz et al. (2025) acknowledge that adding tags improves results on one
dataset, no comprehensive comparison of description strategies across multiple
datasets has been conducted. This gap motivates the need for a systematic
investigation of item description strategies in LLM-based sequential recommendation.


% ── Section 3: Main Contributions ──────────────────────────
\section*{3. Main Contributions}

This work presents two main contributions:

\medskip
\noindent\textbf{Contribution 1: Length-Adaptive Hybrid GNN-Transformer Model.}
We propose a hybrid model that combines Graph Neural Networks and Transformer-based
sequential encoders through an adaptive fusion mechanism. Unlike existing approaches
that apply a uniform fusion strategy for all users, our model dynamically adjusts
the contribution of each component based on the length of the user's interaction
history. This allows the model to better leverage global collaborative signals for
sparse users while exploiting complex sequential patterns for users with rich
histories.

\medskip
\noindent\textbf{Contribution 2: Systematic Study of Item Description Strategies
for LLM-Based Sequential Recommendation.}
We conduct the first systematic investigation of how the formulation of item textual
descriptions affects the quality of LLM-generated embeddings in the context of
sequential recommendation. We compare six description strategies across three
benchmark datasets, providing practical guidance for designing effective
LLM-based recommendation pipelines.


% ── Section 4: Memoir Organization ─────────────────────────
\section*{4. Memoir Organization}

The remainder of this memoir is organized as follows:

\medskip
\textbf{Chapter 1} presents the foundations and state of the art in sequential
recommendation systems, covering classical approaches, sequence-based and
graph-based deep learning models, the integration of Large Language Models, and
evaluation metrics. The chapter concludes by identifying the research gaps that
motivate our contributions.

\textbf{Chapter 2} introduces our first contribution: the Length-Adaptive Hybrid
GNN-Transformer model. We describe the limitations of existing models, formalize
the problem, and present the proposed architecture and fusion mechanism in detail.

\textbf{Chapter 3} reports the experimental evaluation of Contribution 1. We
describe the datasets, preprocessing protocol, baselines, hyperparameter settings,
and discuss the results including an ablation study.

\textbf{Chapter 4} presents our second contribution: a systematic study of item
description strategies for LLM-based sequential recommendation. We introduce the
LLM2Sequential pipeline and describe the four families of description strategies
evaluated in this work.

\textbf{Chapter 5} reports the experimental evaluation of Contribution 2, including
a comparison of all description strategies and a benchmark against the reference
work of Boz et al. (2025).

Finally, a \textbf{General Conclusion} summarizes the key findings of this work,
discusses its limitations, and outlines promising directions for future research.


% ============================================================
% CHAPTER 1 — FOUNDATIONS AND STATE OF THE ART
% ============================================================
\chapter{Foundations and State of the Art}
\label{ch:state-of-the-art}

% ============================================================
\section{Sequential Recommendation}
% ============================================================

\subsection{From Static to Sequential Models}

The classical recommendation approaches described in the previous
section --- collaborative filtering, content-based filtering, and
their hybrid combinations --- treat a user's interaction history as
an unordered set of signals. They aggregate all past interactions
into a single static user representation, ignoring the temporal
structure that underlies real-world behavior. In practice, however,
user preferences are not static: a user who successively interacted
with an action film, a thriller, and a crime drama has revealed a
meaningful and evolving intent that a bag-of-interactions model
cannot capture. \textcite{wang2019sequential} introduced the
sequential recommendation paradigm to address this limitation,
explicitly modeling the \textbf{temporal order} of user interactions
to predict the next item a user is most likely to engage with.

Sequential recommendation has become one of the most active research
directions in the field, driven by the recognition that the order of
interactions encodes rich signals about dynamic user preferences,
contextual shifts, and short-term intent. Unlike traditional
collaborative filtering, which aims to capture long-term, stable
preferences, sequential models are designed to detect \textit{when}
and \textit{why} a user's interest shifts, enabling finer-grained
and more timely recommendations \parencite{wang2019sequential}.

% ------------------------------------------------------------
\subsection{Problem Formulation and Notation}
% ------------------------------------------------------------

Let $\mathcal{U} = \{u_1, u_2, \ldots, u_m\}$ denote the set of
users and $\mathcal{I} = \{v_1, v_2, \ldots, v_n\}$ denote the set
of items, where $m = |\mathcal{U}|$ and $n = |\mathcal{I}|$.
For each user $u \in \mathcal{U}$, we define their \textbf{interaction
sequence} as the chronologically ordered list of items they have
interacted with:
\begin{equation}
  S_u = \left( v_1^u,\, v_2^u,\, \ldots,\, v_{T_u}^u \right),
\end{equation}
where $v_t^u \in \mathcal{I}$ denotes the item interacted with at
time step $t$, and $T_u = |S_u|$ is the total length of the sequence
for user $u$. Interactions can represent clicks, purchases, views,
or ratings depending on the application domain.

Each item $v_i \in \mathcal{I}$ is associated with a learnable
embedding vector $\mathbf{e}_i \in \mathbb{R}^d$, where $d$ is the
latent dimension. The full item embedding matrix is denoted
$\mathbf{E} \in \mathbb{R}^{n \times d}$. Given the interaction
sequence $S_u$, the corresponding sequence of item embeddings forms
the input matrix:
\begin{equation}
  \mathbf{E}_u^{(t)} =
    \left[ \mathbf{e}_{v_1^u},\, \mathbf{e}_{v_2^u},\,
           \ldots,\, \mathbf{e}_{v_t^u} \right]
    \in \mathbb{R}^{t \times d}.
\end{equation}

\subsection{The Next-Item Prediction Task}

The central task of sequential recommendation is
\textbf{next-item prediction}: given the subsequence of interactions
observed up to time step $t$, denoted
$S_u^{(t)} = (v_1^u, \ldots, v_t^u)$, the model must predict the
item $v_{t+1}^u$ that the user will interact with at the next time
step. As formalized by \textcite{kang2018self}, the model learns a
scoring function $f_\theta$ parameterized by $\theta$ such that:
\begin{equation}
  \hat{v}_{t+1}^u = \underset{v \in \mathcal{I} \setminus S_u^{(t)}}
                    {\arg\max} \; f_\theta\!\left(S_u^{(t)},\, v\right),
\end{equation}
where $\mathcal{I} \setminus S_u^{(t)}$ denotes the set of items not
yet seen by user $u$ up to time $t$. The function $f_\theta$ maps
the observed sequence and a candidate item to a relevance score,
and the recommendation is the ranked list of top-$K$ items with
the highest scores.

\subsection{Training Objective}

Sequential recommendation models are typically trained using a
\textbf{next-item prediction objective} over all time steps in the
training sequences. The most common formulation uses a
\textbf{cross-entropy loss} over the full item vocabulary:
\begin{equation}
  \mathcal{L} = - \sum_{u \in \mathcal{U}} \sum_{t=1}^{T_u - 1}
    \log P\!\left( v_{t+1}^u \;\middle|\; S_u^{(t)};\, \theta \right),
\end{equation}
where $P(v_{t+1}^u \mid S_u^{(t)};\, \theta)$ is the model's
predicted probability for the ground-truth next item given the
observed prefix sequence. An alternative is the
\textbf{Bayesian Personalized Ranking (BPR) loss}
\parencite{hidasi2016gru4rec}, which optimizes pairwise preferences
between the ground-truth item and a randomly sampled negative item,
and is preferred when full softmax over the item vocabulary is
computationally prohibitive.

\subsection{Data Splits and Evaluation Protocol}

Following the standard leave-one-out evaluation protocol established
by \textcite{kang2018self}, the interaction sequence of each user is
split chronologically: the \textbf{last item} of the sequence is held
out as the \textbf{test} target, the \textbf{second-to-last item} as
the \textbf{validation} target, and all preceding items form the
\textbf{training} sequence. This strategy ensures that the model is
always evaluated on its ability to predict a truly unseen future
interaction, strictly respecting the temporal ordering of the data
\parencite{wang2019sequential}.

The models reviewed in the following sections propose different
architectures to instantiate the scoring function $f_\theta$,
each making specific assumptions about which aspects of the
interaction sequence are most informative --- and each exhibiting
distinct limitations that our contributions address.

% ============================================================
\section{RNN-Based Sequential Models: GRU4Rec}
% ============================================================

\subsection{Motivation and Core Idea}

Prior to the introduction of neural sequence models, session-based
recommendation relied primarily on memory-based approaches such as
item-to-item similarity, which model only the most recent interaction
and fail to capture the cumulative context of an entire session.
\textcite{hidasi2016gru4rec} addressed this limitation by proposing
\textbf{GRU4Rec}, the first model to apply Gated Recurrent Units
(GRUs) to the recommendation domain. The central idea is to treat
the user's interaction history as a time-ordered sequence and process
it recurrently, allowing the model to maintain a hidden state that
accumulates contextual information across all past interactions in
the session.

\subsection{Architecture}

GRU4Rec processes the sequence of interacted items through one or
more stacked GRU layers \parencite{hidasi2016gru4rec}. At each time
step $t$, the item $v_t$ is mapped to its embedding vector
$\mathbf{e}_t \in \mathbb{R}^d$, which serves as the input to the
recurrent layer. The GRU updates its hidden state $\mathbf{h}_t$
according to the standard gating mechanism:
\begin{equation}
  \mathbf{h}_t = (1 - \mathbf{z}_t) \odot \mathbf{h}_{t-1}
               + \mathbf{z}_t \odot \tilde{\mathbf{h}}_t,
\end{equation}
where $\mathbf{z}_t$ is the \textit{update gate} that controls how
much of the previous hidden state is retained, $\tilde{\mathbf{h}}_t$
is the candidate hidden state computed from the current input and the
reset gate, and $\odot$ denotes element-wise multiplication. The
hidden state $\mathbf{h}_t$ at the final time step encodes the
cumulative sequential context of the session and is used to score
all candidate items via a dot product with their embeddings.

A key contribution of \textcite{hidasi2016gru4rec} is the design of
a \textbf{session-parallel mini-batch} training strategy, which
simultaneously processes multiple user sessions to fully exploit GPU
parallelism. The model is trained with a ranking loss --- either
cross-entropy or Bayesian Personalized Ranking (BPR) --- specifically
adapted to optimize top-$K$ recommendation quality rather than
generic next-token prediction accuracy.

\subsection{Limitations}

Despite its pioneering contribution, GRU4Rec suffers from a
fundamental constraint inherited from all recurrent architectures:
the \textbf{vanishing gradient problem}. As the sequence grows
longer, gradients propagated back through time diminish
exponentially, making it increasingly difficult for the model to
learn dependencies between items that are far apart in the
interaction history \parencite{hidasi2016gru4rec}. Consequently,
GRU4Rec is effective at capturing short-term, local patterns within
a session but struggles to model long-range dependencies that span
many interactions. Furthermore, the strictly sequential left-to-right
processing prevents the model from attending selectively to the most
relevant past interactions, and the absence of any graph-based
component means that global collaborative signals --- relationships
between items derived from the full user-item interaction graph ---
are entirely ignored. These limitations directly motivate the
attention-based and graph-enhanced approaches reviewed in the
following sections.

% ============================================================
\section{CNN-Based Sequential Models: Caser}
% ============================================================

\subsection{Motivation and Core Idea}

While GRU4Rec demonstrated the value of modeling sequential order
through recurrent processing, it captures patterns at the
\textit{point level} --- each prediction depends primarily on the
single most recent hidden state. \textcite{tang2018caser} identified
a complementary type of pattern that recurrent models fundamentally
miss: \textbf{union-level sequential patterns}, where the
\textit{combination} of several recent items jointly predicts the
next interaction, rather than any single item in isolation. To
capture these patterns, \textcite{tang2018caser} proposed
\textbf{Caser} (Convolutional Sequence Embedding Recommendation
Model), which reframes the sequential recommendation problem as a
2D image recognition task.

\subsection{Architecture}

The core idea of Caser is to embed the $L$ most recently interacted
items into a 2D matrix --- an ``image'' in the time and latent
spaces --- and apply convolutional filters to extract local sequential
patterns \parencite{tang2018caser}. Formally, given the $L$ most
recent items in user $u$'s interaction history, their embeddings are
stacked row-wise to form the input matrix
$\mathbf{E}^{(L)} \in \mathbb{R}^{L \times d}$,
where each row $\mathbf{e}_{v_t}$ corresponds to the embedding of
the item at position $t$.

Two types of convolutional filters are then applied over this matrix:
\begin{itemize}
  \item \textbf{Horizontal filters} of shape $h \times d$ (with
        $h \in \{1, 2, \ldots, L\}$) slide across consecutive rows
        and capture \textit{union-level} patterns spanning $h$ recent
        interactions.
  \item \textbf{Vertical filters} of shape $L \times 1$ slide across
        each latent dimension column independently and capture
        \textit{point-level} patterns across all $L$ positions.
\end{itemize}

The outputs of all convolutional filters are concatenated and merged
with a \textbf{user embedding} $\mathbf{p}_u \in \mathbb{R}^d$,
which encodes the user's long-term general preferences independently
of the recent sequence. The combined representation is passed through
a fully connected layer followed by a softmax to produce the ranking
scores over all candidate items \parencite{tang2018caser}.

\subsection{Limitations}

Despite its ability to capture union-level patterns that elude
recurrent models, Caser's design imposes a \textbf{fixed receptive
field}: only the $L$ most recent interactions are encoded in the
input matrix, and all interactions beyond this window are entirely
discarded \parencite{tang2018caser}. This hard truncation makes the
model structurally incapable of capturing long-range dependencies
that span more than $L$ steps. The choice of $L$ is a critical
hyperparameter --- a small $L$ limits the model's contextual horizon,
while a large $L$ increases computational cost and introduces noise
from older, potentially irrelevant interactions. Moreover, like
GRU4Rec, Caser operates exclusively on the local interaction
sequence of each user and contains no mechanism to incorporate
global collaborative signals from the broader user-item interaction
graph, leaving a significant source of information unexploited.
These limitations motivate the transition toward attention-based
architectures capable of selectively attending to any position in
the full interaction history.

% ============================================================
\section{Transformer-Based Sequential Models: SASRec}
% ============================================================

\subsection{Motivation and Core Idea}

Both GRU4Rec and Caser suffer from a shared structural limitation:
they either process sequences strictly left-to-right with a fixed
hidden state (RNN), or apply convolution over a fixed-size window
(CNN), making it impossible to selectively focus on the most
relevant items across the \textit{entire} interaction history.
\textcite{kang2018self} proposed \textbf{SASRec} (Self-Attentive
Sequential Recommendation) to overcome this limitation by applying
the \textbf{self-attention mechanism} --- the core building block of
the Transformer architecture --- directly to sequential
recommendation. The key insight is that at each time step, the model
should be able to attend to any subset of past interactions, weighting
them according to their relevance to the current prediction, rather
than relying on a fixed recurrence or a local convolutional window.

\subsection{Architecture}

SASRec operates on the sequence of the $n$ most recent items in the
user's history \parencite{kang2018self}. Each item $v_t$ is
represented by a learnable embedding $\mathbf{e}_t \in \mathbb{R}^d$,
to which a \textbf{learnable positional embedding} $\mathbf{p}_t
\in \mathbb{R}^d$ is added to preserve the temporal ordering of the
sequence:
\begin{equation}
  \hat{\mathbf{e}}_t = \mathbf{e}_t + \mathbf{p}_t.
\end{equation}

The resulting sequence of embeddings is passed through a stack of
$B$ identical \textbf{self-attention blocks}. Within each block, the
scaled dot-product attention mechanism computes, for each position
$t$, a weighted sum over all preceding positions:
\begin{equation}
  \text{Attention}(\mathbf{Q}, \mathbf{K}, \mathbf{V}) =
    \text{softmax}\!\left(
      \frac{\mathbf{Q}\mathbf{K}^\top}{\sqrt{d}}
      + \mathbf{M}
    \right)\mathbf{V},
\end{equation}
where $\mathbf{Q}$, $\mathbf{K}$, and $\mathbf{V}$ are the query,
key, and value matrices derived by linear projections of the input,
$d$ is the embedding dimension, and $\mathbf{M}$ is a
\textbf{causal mask} that sets future positions to $-\infty$ before
the softmax, ensuring that the prediction at position $t$ can only
attend to items at positions $1, \ldots, t-1$
\parencite{kang2018self}. A point-wise feed-forward network and
layer normalization are applied after each attention block. The final
hidden representation at the last position is used to rank all
candidate items via a dot product with their embeddings.

\subsection{Limitations}

SASRec represents a major advance over RNN and CNN-based models by
enabling long-range, selective attention across the full interaction
history. However, it presents two notable limitations. First, the
\textbf{unidirectional causal mask} restricts each position to
attending only to past items, which prevents the model from
leveraging bidirectional context that could further enrich item
representations --- a gap later addressed by BERT4Rec
\parencite{sun2019bert4rec}. Second, and more critically for our
work, SASRec processes each user's interaction sequence
\textbf{in complete isolation}: it has no access to the global
user-item interaction graph, and therefore cannot exploit the rich
collaborative signals that emerge from the co-interaction patterns
of all users across the entire dataset \parencite{kang2018self}.
Two users who have interacted with many of the same items receive
no shared signal within the SASRec framework, as item embeddings
are learned solely from individual sequences. This absence of global
collaborative information directly motivates the integration of
Graph Neural Networks into sequential recommendation pipelines,
as reviewed in Section~\ref{sec:lightgcn}.


% ============================================================
\section{Bidirectional Transformer Models: BERT4Rec}
% ============================================================

\subsection{Motivation and Core Idea}

SASRec demonstrated that self-attention is highly effective for
sequential recommendation, yet its \textbf{unidirectional causal
architecture} processes each position by attending only to past
items. \textcite{sun2019bert4rec} argued that this left-to-right
constraint artificially limits the representational power of item
embeddings within the sequence: in practice, the context
surrounding an item --- both before \textit{and} after it ---
can be informative for learning richer representations.
Inspired by the success of BERT in natural language processing,
\textcite{sun2019bert4rec} proposed \textbf{BERT4Rec}, which
applies \textbf{deep bidirectional self-attention} to model user
interaction sequences, allowing every item to attend to all other
items in the sequence regardless of their position.

\subsection{Architecture}

BERT4Rec adopts the same multi-layer bidirectional Transformer
encoder architecture as BERT \parencite{sun2019bert4rec}. The
interaction sequence of length $n$ is embedded into a matrix, to
which learnable positional embeddings are added. The model then
applies $L$ stacked bidirectional self-attention blocks, where
each position can attend to \textit{all} other positions in the
sequence --- both left and right context --- without any causal
mask.

The key challenge of training a bidirectional model for
recommendation is \textbf{information leakage}: if the model can
see all positions simultaneously, predicting the next item becomes
trivial. \textcite{sun2019bert4rec} addressed this by introducing
a \textbf{Cloze task} (also known as Masked Language Modeling)
adapted for recommendation. During training, a random fraction
$\rho$ of items in the sequence are replaced by a special
\texttt{[MASK]} token, and the model is trained to predict the
original items at the masked positions by conditioning jointly on
both left and right context:
\begin{equation}
  \mathcal{L}_{\text{Cloze}} =
    -\sum_{v_m \in S_{\text{masked}}}
      \log P\!\left( v_m = v_m^* \;\middle|\; \tilde{S}_u \right),
\end{equation}
where $S_{\text{masked}}$ is the set of masked positions,
$v_m^*$ is the true item at masked position $m$, and
$\tilde{S}_u$ is the corrupted sequence with masked tokens
\parencite{sun2019bert4rec}. This formulation generates multiple
training samples per sequence (one per masked position), producing
a richer training signal than the single next-item prediction
objective of SASRec. At inference time, the last item of the
sequence is replaced by \texttt{[MASK]}, and the model's output
at that position is used to rank all candidate items.

\subsection{Limitations}

BERT4Rec consistently outperforms SASRec and earlier sequential
models on standard benchmarks, yet it introduces two notable
limitations. First, the \textbf{train-test discrepancy}: the model
is trained to predict masked items at arbitrary positions, but is
evaluated exclusively on the last position, creating a mismatch
between the training objective and the inference task
\parencite{sun2019bert4rec}. Second, and most critically for this
work, BERT4Rec --- like all models reviewed so far --- processes
each user's sequence \textbf{in complete isolation} from the rest
of the dataset. It has no mechanism to leverage the global
structure of the user-item interaction graph, and therefore cannot
exploit the collaborative signals that emerge from the collective
behavior of all users. Item representations are learned purely
from individual interaction sequences, with no awareness of which
items are frequently co-interacted with across the population.
This fundamental gap motivates the graph-based approach reviewed
in the following section.

% ============================================================
\section{Graph Neural Network Models: LightGCN}
\label{sec:lightgcn}
% ============================================================

\subsection{Motivation and Core Idea}

The sequential models reviewed so far --- GRU4Rec, Caser, SASRec,
and BERT4Rec --- all share a common blind spot: they model each
user's interaction history in isolation, with no access to the
global structure of the user-item interaction graph. This means
they cannot exploit \textbf{high-order collaborative signals} ---
the rich patterns that emerge when items are co-interacted with by
many users across the entire dataset. Graph Convolution Network
(GCN)-based approaches address this limitation by propagating
information across the interaction graph. However, \textcite{he2020lightgcn}
demonstrated through systematic ablation experiments that the two
most common operations in standard GCNs --- \textbf{feature
transformation} and \textbf{nonlinear activation} --- contribute
negligibly to recommendation performance and even harm it by
increasing training difficulty. They proposed \textbf{LightGCN},
which strips the GCN to its single essential operation ---
\textbf{neighborhood aggregation} --- and achieves approximately
16\% relative improvement over the prior state-of-the-art NGCF
under identical experimental conditions \parencite{he2020lightgcn}.

\subsection{Architecture}

LightGCN constructs a bipartite user-item interaction graph
$\mathcal{G} = (\mathcal{U} \cup \mathcal{I},\, \mathcal{E})$,
where each edge $(u, i) \in \mathcal{E}$ indicates that user $u$
has interacted with item $i$ \parencite{he2020lightgcn}. Starting
from learnable initial embeddings $\mathbf{e}_u^{(0)}$ and
$\mathbf{e}_i^{(0)}$, LightGCN performs $K$ rounds of
\textbf{linear neighborhood aggregation}. At each layer $k$, the
embedding of a node is updated as the normalized weighted sum of
its neighbors' embeddings from the previous layer:
\begin{equation}
  \mathbf{e}_u^{(k)} =
    \sum_{i \in \mathcal{N}_u}
      \frac{1}{\sqrt{|\mathcal{N}_u|}\sqrt{|\mathcal{N}_i|}}
      \mathbf{e}_i^{(k-1)},
  \qquad
  \mathbf{e}_i^{(k)} =
    \sum_{u \in \mathcal{N}_i}
      \frac{1}{\sqrt{|\mathcal{N}_i|}\sqrt{|\mathcal{N}_u|}}
      \mathbf{e}_u^{(k-1)},
\end{equation}
where $\mathcal{N}_u$ and $\mathcal{N}_i$ denote the sets of
items interacted with by user $u$ and users who interacted with
item $i$, respectively. There are \textbf{no weight matrices} and
\textbf{no activation functions} --- the propagation is purely
linear \parencite{he2020lightgcn}.

The final user and item embeddings are computed as the
\textbf{layer-combination} of all $K$ propagation layers:
\begin{equation}
  \mathbf{e}_u = \sum_{k=0}^{K} \alpha_k\, \mathbf{e}_u^{(k)},
  \qquad
  \mathbf{e}_i = \sum_{k=0}^{K} \alpha_k\, \mathbf{e}_i^{(k)},
\end{equation}
where $\alpha_k = \frac{1}{K+1}$ is a uniform layer weight.
This multi-layer combination captures embeddings that encode
increasingly high-order neighborhood structure: layer 1 captures
direct neighbors, layer 2 captures two-hop neighbors, and so on.
The predicted relevance score for a user-item pair is then computed
as the inner product $\hat{r}_{ui} = \mathbf{e}_u^\top \mathbf{e}_i$,
and the model is trained with the BPR loss over observed and
sampled negative interactions \parencite{he2020lightgcn}.

\subsection{Limitations}

Despite its effectiveness at capturing global collaborative signals
through graph propagation, LightGCN has a fundamental limitation
from the perspective of sequential recommendation: it is entirely
\textbf{order-agnostic}. The bipartite interaction graph contains
no temporal information --- all historical interactions are treated
as equally weighted, static evidence, regardless of when they
occurred \parencite{he2020lightgcn}. LightGCN cannot distinguish
a user who watched a film last year from one who watched it
yesterday, nor can it detect any drift in user preferences over
time. As a result, it fails to capture the dynamic, short-term
intent that is central to the sequential recommendation task.
This complementarity --- sequential models capturing temporal
dynamics but lacking global signals, and LightGCN capturing global
signals but lacking temporal dynamics --- is precisely the gap that
the hybrid architecture proposed in this work seeks to bridge.

% ============================================================
\section{Large Language Models for Sequential Recommendation}
% ============================================================

\subsection{From ID-Based to Semantic Item Representations}

All models reviewed in the previous sections --- GRU4Rec, Caser,
SASRec, BERT4Rec, and LightGCN --- represent items as
\textbf{anonymous learnable ID embeddings}: each item is assigned
a randomly initialized vector that is optimized purely from the
observed interaction data. While effective when sufficient
interaction history is available, this paradigm has a fundamental
limitation: it treats items as opaque identifiers with no intrinsic
meaning, discarding the rich semantic content embedded in item
attributes such as titles, genres, descriptions, and metadata.
This makes such models particularly vulnerable to the
\textbf{cold-start problem}, where items with few or no interactions
have poorly learned embeddings regardless of how descriptive their
content is.

The emergence of \textbf{Large Language Models (LLMs)} --- deep
neural networks pretrained on massive textual corpora via
self-supervised objectives --- has introduced a fundamentally
different source of item representations. Models such as GPT,
LLaMA, and BERT encode rich semantic and world knowledge acquired
during pretraining, enabling them to produce high-quality
embeddings from textual item descriptions even in the complete
absence of interaction data \parencite{wu2023survey}.

\subsection{The LLM2Sequential Framework}

\textcite{boz2025improving} conducted a comprehensive empirical
study of three orthogonal strategies for integrating LLMs into
sequential recommendation pipelines, evaluated across three
datasets and multiple baseline models. Their most impactful
approach, which they call \textbf{LLM2Sequential}, consists of
replacing the randomly initialized item embedding matrix of an
existing sequential model with embeddings pre-computed by an LLM
from item textual descriptions. These LLM-generated embeddings
are obtained by passing each item's description through a
pretrained language model and extracting the resulting dense
representation. The embedding matrix is then used to initialize
--- and optionally fine-tune --- the item embeddings of models
such as BERT4Rec or SASRec.

\textcite{boz2025improving} demonstrate that this simple
initialization strategy yields substantial performance gains:
initializing BERT4Rec with LLM embeddings improves NDCG by
15--20\% compared to the standard randomly initialized BERT4Rec
on their benchmark datasets. These results establish that
\textbf{semantic item embeddings from LLMs are highly
complementary to the sequential modeling capabilities of
Transformer-based architectures}, providing a strong initialization
that accelerates convergence and enriches item representations
with content-level signals.

\subsection{The Open Question: Description Strategy}

Despite the strong results reported by \textcite{boz2025improving},
their study leaves a critical question unaddressed: \textbf{how
should an item be described to the LLM in order to maximize the
quality of the resulting embedding?} In their experiments, items
are described primarily by their title, with one dataset including
additional tags. No systematic comparison of different description
strategies is conducted across datasets.

This gap is significant because the textual input provided to the
LLM fundamentally determines what semantic information is encoded
in the resulting embedding. A title-only description captures
the most basic item identity, while richer formulations --- such
as metadata-enriched profiles combining genre, director, cast and
synopsis, or user-centric narratives summarizing community
preferences --- may provide the LLM with substantially more
discriminative information. \textcite{boz2025improving} themselves
acknowledge that ``adding tags improves results on one dataset'',
implicitly suggesting that description richness matters, yet no
principled investigation is provided.

This unresolved question directly motivates the second contribution
of this work: a \textbf{systematic study of item description
strategies} for LLM-based sequential recommendation, presented in
Chapter~4.

% ============================================================
\section{Evaluation Metrics}
% ============================================================

The performance of sequential recommendation models is evaluated
using ranking metrics that measure how well the ground-truth next
item is ranked among all candidate items. Following the standard
evaluation protocol established in the literature
\parencite{kang2018self}, we report three complementary metrics
at cutoff $K$: \textbf{Hit Rate (HR@K)}, \textbf{Normalized
Discounted Cumulative Gain (NDCG@K)}, and \textbf{Mean Reciprocal
Rank (MRR@K)}. These three metrics are computed over the top-$K$
ranked list returned by the model for each test user, where the
ground truth is the single held-out last item from the user's
sequence.

\subsection{Hit Rate at K (HR@K)}

Hit Rate at $K$ (HR@$K$), also referred to as Recall@$K$ in the
next-item prediction setting, measures the fraction of test users
for whom the ground-truth item appears anywhere within the top-$K$
recommended list \parencite{kang2018self}. Formally:
\begin{equation}
  \text{HR@}K = \frac{1}{|\mathcal{U}|}
    \sum_{u \in \mathcal{U}}
    \mathbf{1}\!\left[ \text{rank}(v_{T_u+1}^u) \leq K \right],
\end{equation}
where $v_{T_u+1}^u$ is the ground-truth next item for user $u$,
$\text{rank}(\cdot)$ denotes its position in the ranked list
returned by the model, and $\mathbf{1}[\cdot]$ is the indicator
function. HR@$K$ takes values in $[0, 1]$, where 1 indicates
that the correct item is in the top-$K$ list for every user.
It is a binary metric --- it rewards any correct recommendation
within the top-$K$ regardless of its exact position --- making
it a measure of \textit{coverage} rather than \textit{ranking
quality}.

\subsection{Normalized Discounted Cumulative Gain (NDCG@K)}

NDCG@$K$ is a position-sensitive metric that penalizes correct
recommendations placed lower in the ranked list. It assigns a
higher reward to models that rank the ground-truth item as close
to position 1 as possible \parencite{wang2019sequential}. The
Discounted Cumulative Gain at $K$ for a single user is:
\begin{equation}
  \text{DCG@}K =
    \frac{1}{\log_2(\text{rank}(v_{T_u+1}^u) + 1)}
    \cdot \mathbf{1}\!\left[ \text{rank}(v_{T_u+1}^u) \leq K \right],
\end{equation}
and the normalized version is obtained by dividing by the Ideal
DCG (IDCG), which equals $1 / \log_2(2) = 1$ in the single
ground-truth item setting. The metric averaged over all users is:
\begin{equation}
  \text{NDCG@}K = \frac{1}{|\mathcal{U}|}
    \sum_{u \in \mathcal{U}}
    \frac{1}{\log_2\!\left(\text{rank}(v_{T_u+1}^u) + 1\right)}
    \cdot \mathbf{1}\!\left[ \text{rank}(v_{T_u+1}^u) \leq K \right].
\end{equation}
NDCG@$K$ is the most widely used metric in sequential
recommendation research, as it simultaneously captures both
whether the correct item is retrieved and how highly it is ranked.

\subsection{Mean Reciprocal Rank (MRR@K)}

Mean Reciprocal Rank at $K$ (MRR@$K$) measures the average of
the reciprocal ranks of the ground-truth item across all test
users \parencite{hidasi2016gru4rec}. Formally:
\begin{equation}
  \text{MRR@}K = \frac{1}{|\mathcal{U}|}
    \sum_{u \in \mathcal{U}}
    \frac{1}{\text{rank}(v_{T_u+1}^u)}
    \cdot \mathbf{1}\!\left[ \text{rank}(v_{T_u+1}^u) \leq K \right].
\end{equation}
MRR gives the highest reward when the correct item is ranked first
(reciprocal rank $= 1$) and rapidly decreasing rewards for lower
positions. It is particularly sensitive to the quality of the
top-1 prediction and complements NDCG by emphasizing the
\textit{precision} of the highest-ranked recommendation.

\subsection{Summary and Relation to This Work}

\begin{table}[h]
\centering
\caption{Properties of the three evaluation metrics used in this work}
\small
\begin{tabular}{lccc}
\toprule
\textbf{Property} & \textbf{HR@K} & \textbf{NDCG@K} & \textbf{MRR@K} \\
\midrule
Position-sensitive        & $\times$ & \checkmark & \checkmark \\
Binary (hit/miss)         & \checkmark & $\times$ & $\times$ \\
Rewards rank 1 strongly   & $\times$ & Partial & \checkmark \\
Standard in literature    & \checkmark & \checkmark & \checkmark \\
\bottomrule
\end{tabular}
\end{table}

In this work, all models are evaluated using HR@10, NDCG@10, and
MRR@10, following the standard cutoff value $K = 10$ adopted in
the sequential recommendation literature
\parencite{kang2018self, sun2019bert4rec}. The three metrics
together provide a comprehensive picture of model performance:
HR@10 measures coverage, NDCG@10 measures ranking quality, and
MRR@10 measures top-1 precision.

% ============================================================
\section{Research Gaps and Contributions}
% ============================================================

\subsection{Synthesis of Limitations}

The preceding sections have reviewed the five baseline models that
form the foundation of this work: GRU4Rec, Caser, SASRec, BERT4Rec,
and LightGCN. While each model represents a significant advance in
its respective paradigm, a systematic analysis reveals two
persistent gaps that none of them address.

\begin{table}[ht]
\centering
\caption{Limitations of baseline sequential recommendation models}
\renewcommand{\arraystretch}{1.3}
\footnotesize
\begin{tabularx}{\textwidth}{@{} X c c c c @{}}
\toprule
\textbf{Model} & \textbf{Long-range} & \textbf{Global graph} &
\textbf{Adaptive fusion} & \textbf{Semantic items} \\
\midrule
GRU4Rec \parencite{hidasi2016gru4rec}
  & $\times$ & $\times$ & $\times$ & $\times$ \\
Caser \parencite{tang2018caser}
  & $\times$ & $\times$ & $\times$ & $\times$ \\
SASRec \parencite{kang2018self}
  & \checkmark & $\times$ & $\times$ & $\times$ \\
BERT4Rec \parencite{sun2019bert4rec}
  & \checkmark & $\times$ & $\times$ & $\times$ \\
LightGCN \parencite{he2020lightgcn}
  & $\times$ & \checkmark & $\times$ & $\times$ \\
\midrule
\textbf{This work}
  & \checkmark & \checkmark & \checkmark & \checkmark \\
\bottomrule
\end{tabularx}
\end{table}

\subsection{Gap 1 — Static Fusion of Sequential and Graph Signals}

GRU4Rec and Caser cannot capture long-range dependencies due to
their architectural constraints --- vanishing gradients and fixed
convolutional windows respectively. SASRec and BERT4Rec overcome
this limitation through self-attention, but they process each
user's interaction history in complete isolation from the rest of
the dataset, with no access to the global collaborative signals
encoded in the user-item interaction graph. LightGCN, conversely,
captures these global signals effectively through graph
propagation, but ignores the temporal order of interactions
entirely, treating all historical interactions as equally weighted
static evidence.

Existing hybrid approaches that combine sequential and graph-based
components \parencite{he2020lightgcn, kang2018self} address this
complementarity, but rely on \textbf{static fusion strategies}:
the contribution of each component is fixed uniformly for all
users. This ignores a fundamental disparity in the data --- user
interaction histories vary significantly in length and density.
Users with sparse histories have insufficient sequential signal
for Transformer-based models to operate reliably, and should
benefit more from the global collaborative signal provided by
the graph component. Conversely, users with long and rich
histories provide sufficient sequential patterns for the
Transformer to capture complex temporal dynamics, reducing their
dependence on the graph.

\noindent\textbf{Gap 1:} \textit{No existing model dynamically
adapts the contribution of global graph signals versus local
sequential signals based on the richness of the user's
interaction history.}

\subsection{Gap 2 — No Systematic Study of LLM Description Strategies}

\textcite{boz2025improving} demonstrated that initializing
sequential models with LLM-generated item embeddings yields
substantial performance gains over ID-based representations.
However, their study uses a single fixed description strategy
--- primarily item titles --- and does not systematically
investigate whether the \textit{formulation} of the textual
description provided to the LLM affects the quality of the
resulting embeddings. They acknowledge in their results that
adding tags improves performance on one dataset, implicitly
suggesting that description richness matters, yet provide
no principled analysis.

This gap is significant: the textual input to the LLM is the
sole determinant of what semantic information is encoded in
the embedding. Different formulations --- a bare title, a
metadata-enriched profile, a user-centric narrative, or a
structured hybrid --- capture fundamentally different facets
of the item and may produce embeddings of substantially
different discriminative quality for the recommendation task.

\noindent\textbf{Gap 2:} \textit{No systematic study has
compared the impact of different item description strategies
on LLM embedding quality in the context of sequential
recommendation.}

\subsection{Contributions of This Work}

The two gaps identified above directly motivate the two
contributions of this memoir:

\medskip
\noindent\textbf{Contribution 1 --- Length-Adaptive Hybrid
GNN-Transformer Model.}
We propose a hybrid sequential recommendation model that
combines a LightGCN-based graph encoder and a BERT4Rec-based
Transformer encoder through a \textbf{length-adaptive fusion
mechanism}. Unlike static fusion, our model computes a
user-specific fusion weight $\alpha(u) \in [0, 1]$ as a
monotonically increasing function of the user's interaction
history length $T_u$, dynamically shifting the balance from
graph-dominated representations for sparse users toward
sequence-dominated representations for dense users. This
contribution is presented in Chapter~2 and evaluated in
Chapter~3.

\medskip
\noindent\textbf{Contribution 2 --- Systematic Study of Item
Description Strategies for LLM-Based Sequential Recommendation.}
We conduct the first systematic investigation of how the
formulation of item textual descriptions affects the quality
of LLM-generated embeddings in the LLM2Sequential framework
of \textcite{boz2025improving}. We design and compare six
description strategies across three benchmark datasets, using
BERT4Rec as the downstream sequential model, and provide
quantitative and qualitative guidance for practitioners
designing LLM-augmented recommendation pipelines. This
contribution is presented in Chapter~4 and evaluated in
Chapter~5.



% ============================================================
\chapter{Proposed Methodology: Length-Adaptive Hybrid GNN-Transformer}
\addcontentsline{toc}{chapter}{Chapter 2 — Proposed Methodology}
% ============================================================

\section{Introduction and Motivation}

The state of the art reviewed in Chapter~1 reveals a fundamental
tension at the heart of sequential recommendation research. On one
side, Transformer-based models such as SASRec and BERT4Rec excel at
capturing fine-grained temporal dynamics within individual user
interaction sequences, but they process each user's history in
complete isolation from the rest of the dataset, ignoring the global
co-occurrence structure that connects items across all users. On the
other side, graph-based models such as LightGCN capture precisely
these global collaborative signals through neighborhood propagation
over the user-item interaction graph, but they treat all historical
interactions as an unordered, time-invariant set, entirely discarding
the sequential dynamics that drive next-item prediction.

Hybrid approaches that combine graph and sequential encoders have
been proposed to bridge this gap \parencite{he2020lightgcn,
kang2018self}. However, as identified in Section~1.10, these
approaches suffer from a critical design flaw: they apply a
\textbf{fixed fusion weight} uniformly to all users, regardless
of how much sequential evidence is available for each individual.
This assumption is fundamentally incorrect. The relative reliability
of graph signals versus sequence signals is not constant --- it
varies dramatically with the length of the user's interaction
history. A user with only three recorded interactions provides
insufficient sequential context for self-attention to generalize
reliably; assigning such a user the same fusion weight as a user
with two hundred interactions wastes the compensatory power of the
global graph signal precisely where it is most needed
\parencite{ouledmeriem2025hybrid}.

This observation motivates a principled, user-aware fusion
strategy. Users with \textbf{short interaction histories} should
rely more heavily on the GNN's global co-occurrence signal, which
aggregates information across the full user population and is
independent of sequence length. Users with \textbf{long interaction
histories} provide sufficient sequential evidence for the
Transformer to exploit complex temporal patterns and evolving
preferences, reducing their dependence on the graph component.
The transition between these two regimes should be explicitly
encoded in the model's design rather than left to a single
globally learned scalar \parencite{ouledmeriem2025hybrid}.

This chapter presents \textbf{Contribution 1} of this memoir: the
\textbf{Length-Adaptive Hybrid GNN-Transformer}, a family of four
hybrid architectures that share BERT4Rec as the sequential backbone
and differ in how the complementary signal is constructed and fused.
The unifying insight across all four variants is that GNN signals
rescue sparse and cold-start users where sequential evidence is
insufficient, while the Transformer dominates for active users
with rich interaction histories. The chapter is organized as
follows: Section~2.2 formalizes the problem, Section~2.3 describes
the BERT4Rec backbone, Section~2.4 presents the item co-occurrence
graph and GCN encoder, Section~2.5 introduces the four hybrid
variants in detail, and Sections~2.6--2.7 cover the training
objective and implementation details.


% ============================================================
\section{Problem Formulation}
% ============================================================

\subsection{Basic Notation}

Let $\mathcal{U} = \{u_1, u_2, \ldots, u_m\}$ and
$\mathcal{I} = \{i_1, i_2, \ldots, i_n\}$ denote the sets of
users and items respectively, where $m = |\mathcal{U}|$ and
$n = |\mathcal{I}|$. For each user $u \in \mathcal{U}$, let
\begin{equation}
  S_u = [i_1, i_2, \ldots, i_L]
\end{equation}
denote their interaction history of length $L = T_u$, where each
$i_t \in \mathcal{I}$ is the item consumed at time step $t$,
ordered chronologically. Each item $i \in \mathcal{I}$ is
associated with a learnable embedding $\mathbf{e}_i \in \mathbb{R}^d$,
stored in the embedding matrix $\mathbf{E} \in \mathbb{R}^{n \times d}$,
where $d$ is the latent dimension.

\subsection{Next-Item Prediction Task}

The task of sequential recommendation is defined as follows:
given the prefix sequence $S_u^{(t)} = [i_1, \ldots, i_t]$
observed up to time step $t$, learn a scoring function
\begin{equation}
  f_\Theta\!\left(u,\, S_u^{(t)},\, j\right) \in \mathbb{R}
\end{equation}
that assigns a relevance score to each candidate item
$j \in \mathcal{I}$, where $\Theta$ denotes all trainable
parameters \parencite{ouledmeriem2025hybrid}. At inference time,
items are ranked by their scores and the top-$K$ list is returned.
Performance is evaluated via HR@$K$, NDCG@$K$, and MRR@$K$,
as defined in Section~1.9.

\subsection{The Adaptive Fusion Problem}

Standard sequential recommendation models learn $f_\Theta$ from
the interaction sequence alone, producing a single user
representation $\mathbf{h}_u \in \mathbb{R}^d$ per user. Hybrid
models augment this with a complementary graph-based signal
$\mathbf{g}'_i \in \mathbb{R}^d$ per item, and combine both
through a fusion weight $\alpha \in [0,1]$:
\begin{equation}
  \mathbf{h}_i^{(u)} = \alpha\, \mathbf{e}_i
                     + (1 - \alpha)\, \mathbf{g}'_i.
\end{equation}

Existing hybrid approaches fix $\alpha$ as a single global scalar
shared across all users \parencite{ouledmeriem2025hybrid}. The
central problem addressed in this chapter is to replace this
static assignment with a \textbf{user-specific fusion weight}
$\alpha(u) \in [0,1]$ that adapts to the length $L(u)$ of each
user's interaction history:
\begin{equation}
  \mathbf{h}_i^{(u)} = \alpha(u)\, \mathbf{e}_i
                     + \bigl(1 - \alpha(u)\bigr)\, \mathbf{g}'_i,
\end{equation}
where $\alpha(u)$ is a monotonically non-decreasing function of
$L(u)$ --- assigning more weight to the sequential embedding
$\mathbf{e}_i$ as interaction history grows, and more weight to
the graph embedding $\mathbf{g}'_i$ as history shrinks.

\subsection{User History Length Groups}

Following \textcite{ouledmeriem2025hybrid}, we partition users
into three groups based on interaction history length, defined
by two thresholds $L^*_{\text{short}} = 10$ and
$L^*_{\text{long}} = 30$:

\begin{equation}
  \alpha(u) =
  \begin{cases}
    \alpha_{\text{short}} & \text{if } L(u) \leq 10 \\
    \alpha_{\text{mid}}   & \text{if } 10 < L(u) \leq 30 \\
    \alpha_{\text{long}}  & \text{if } L(u) > 30
  \end{cases}
\end{equation}

The rationale for each bin is grounded in the complementary
strengths of the two encoders:

\begin{itemize}
  \item \textbf{Short-history users} ($L(u) \leq 10$): self-attention
        lacks sufficient context to generalize reliably from fewer
        than 10 interactions. Assigning a small $\alpha_{\text{short}}$
        places greater weight on the GNN's global co-occurrence signal,
        which aggregates evidence across the full user population
        independently of sequence length.

  \item \textbf{Medium-history users} ($10 < L(u) \leq 30$): both
        signals are partially informative. A balanced
        $\alpha_{\text{mid}}$ allows the model to exploit both
        graph-level structure and emerging sequential patterns.

  \item \textbf{Long-history users} ($L(u) > 30$): sufficient
        sequential evidence is available for the Transformer to
        capture complex temporal dynamics. Assigning a large
        $\alpha_{\text{long}}$ preserves the full expressive power
        of the BERT4Rec sequence embedding.
\end{itemize}

The three values $(\alpha_{\text{short}}, \alpha_{\text{mid}},
\alpha_{\text{long}})$ are treated as hyperparameters tuned on
the validation set. In our experiments, the optimal values found
are $\alpha_{\text{short}} = 0.3$, $\alpha_{\text{mid}} = 0.5$,
and $\alpha_{\text{long}} = 0.7$ \parencite{ouledmeriem2025hybrid}.
This mechanism requires no additional trainable parameters beyond
these three scalars and adds negligible computational overhead.


% ============================================================
\section{BERT4Rec Backbone}
% ============================================================

\subsection{Overview}

All four hybrid variants proposed in this work share
\textbf{BERT4Rec} \parencite{sun2019bert4rec} as their sequential
encoder. BERT4Rec adapts the masked language modeling objective
of BERT to the recommendation setting, applying deep bidirectional
self-attention to model the full context of a user's interaction
sequence. It was selected as the backbone for three reasons: it
consistently achieves the strongest performance among purely
sequential baselines on the MovieLens benchmarks used in this
work \parencite{ouledmeriem2025hybrid}; its bidirectional attention
allows every item in the sequence to attend to all other items,
producing richer contextual representations than unidirectional
models such as SASRec; and its modular architecture makes it
straightforward to inject auxiliary graph-based signals at the
embedding level before the Transformer blocks process the sequence.

\subsection{Input Representation}

Each item $i \in \mathcal{I}$ is assigned a learnable embedding
$\mathbf{e}_i \in \mathbb{R}^d$, stored in the embedding matrix
$\mathbf{E} \in \mathbb{R}^{|\mathcal{I}| \times d}$. For a
sequence of length $L$, a set of learnable positional embeddings
$\mathbf{P} \in \mathbb{R}^{L \times d}$ is added element-wise to
encode the temporal ordering of the interactions. The input
representation at position $t$ is:
\begin{equation}
  \mathbf{H}_t^{(0)} = \mathbf{e}_{i_t} + \mathbf{p}_t,
  \quad t = 1, \ldots, L,
\end{equation}
where $\mathbf{p}_t \in \mathbb{R}^d$ is the positional embedding
at position $t$ \parencite{ouledmeriem2025hybrid}. The full input
matrix $\mathbf{H}^{(0)} \in \mathbb{R}^{L \times d}$ is then
passed through $B$ stacked Transformer blocks.

\subsection{Self-Attention Blocks}

Each of the $B$ Transformer blocks applies \textbf{multi-head
self-attention} followed by a position-wise feed-forward network,
residual connections, and layer normalization
\parencite{vaswani2017attention}. Within each block $\ell$, the
query, key, and value matrices are obtained by linear projections
of the input from the previous block:
\begin{equation}
  \mathbf{Q} = \mathbf{H}^{(\ell-1)} \mathbf{W}^Q, \quad
  \mathbf{K} = \mathbf{H}^{(\ell-1)} \mathbf{W}^K, \quad
  \mathbf{V} = \mathbf{H}^{(\ell-1)} \mathbf{W}^V,
\end{equation}
where $\mathbf{W}^Q, \mathbf{W}^K, \mathbf{W}^V \in
\mathbb{R}^{d \times d_k}$ are learnable projection matrices and
$d_k$ is the key dimension. The scaled dot-product attention is
then computed as:
\begin{equation}
  \text{Attention}(\mathbf{Q}, \mathbf{K}, \mathbf{V}) =
    \text{softmax}\!\left(
      \frac{\mathbf{Q} \mathbf{K}^\top}{\sqrt{d_k}}
    \right) \mathbf{V}.
\end{equation}

Unlike SASRec, which applies a causal mask to enforce
left-to-right processing, BERT4Rec uses \textbf{full bidirectional
attention}: every position can attend to all other positions in
the sequence simultaneously, producing richer contextual item
representations \parencite{sun2019bert4rec}. After $B$ blocks,
the final hidden representation at the last position is taken as
the \textbf{sequence embedding}:
\begin{equation}
  \mathbf{h}_u = \mathbf{H}_L^{(B)} \in \mathbb{R}^d.
\end{equation}

\subsection{Item Scoring and Cloze Training}

The relevance score for a candidate item $j$ given user $u$ is
computed as the inner product between the sequence embedding and
the item embedding:
\begin{equation}
  s_{uj} = \mathbf{h}_u^\top \mathbf{e}_j.
\end{equation}

During training, BERT4Rec applies a \textbf{Cloze task}: a random
fraction $\rho$ of items in the input sequence are replaced by a
special \texttt{[MASK]} token, and the model is trained to
predict the original items at the masked positions by attending
jointly to both left and right context \parencite{sun2019bert4rec}.
This produces multiple training samples per sequence and generates
a richer training signal than the single next-item prediction
objective of SASRec. At inference time, the last item in the
sequence is masked and the model's output at that position is
used to rank all candidate items.

\subsection{Why BERT4Rec as the Backbone}

The choice of BERT4Rec over SASRec as the shared backbone is
motivated by three empirical observations from our experiments
\parencite{ouledmeriem2025hybrid}. First, BERT4Rec achieves
higher HR@10, NDCG@10, and MRR@10 than SASRec on both ML-100K
and ML-1M under identical training conditions. Second, the
bidirectional attention mechanism produces richer item-level
contextual representations that are better suited to receive
injected graph signals at the embedding level. Third, the
masked training objective naturally accommodates the modified
embedding table produced by the fusion layer, since the model
must reconstruct masked items from a globally enriched
representation rather than relying on raw ID embeddings alone.

% ============================================================
\section{Item Co-occurrence Graph and GCN Encoder}
% ============================================================

\subsection{Motivation}

The BERT4Rec backbone described in Section~2.3 encodes each user's
interaction history independently, with no access to the structural
relationships between items that emerge from the collective behavior
of all users. Two items that are frequently co-interacted with
across the entire user population share a strong implicit
relationship that purely sequential models cannot detect. To capture
these \textbf{global structural signals}, we construct an
\textbf{item co-occurrence graph} from all training sequences and
encode it with a Graph Convolutional Network (GCN)
\parencite{kipf2017gcn}, producing graph-enhanced item embeddings
that complement the sequential representations of BERT4Rec.

\subsection{Graph Construction}

Given the set of all training interaction sequences
$\{S_u\}_{u \in \mathcal{U}}$, we construct an undirected weighted
item co-occurrence graph $\mathcal{G} = (\mathcal{I}, \mathcal{E})$,
where nodes are items and edges encode co-occurrence frequency
\parencite{ouledmeriem2025hybrid}. For each user sequence $S_u$,
a sliding window of size $w$ is applied over the sequence. For
every window, an undirected edge is created between each pair of
co-occurring items $(a, b)$ and their co-occurrence count $c_{ab}$
is incremented by one. After aggregating over all user sequences,
edge weights are defined as:
\begin{equation}
  w_{ab} =
  \begin{cases}
    c_{ab} & \text{if } c_{ab} \geq \tau \\
    0      & \text{otherwise,}
  \end{cases}
\end{equation}
where $\tau$ is a minimum co-occurrence threshold that prunes
rare and noisy edges. In our experiments, we use a window size
$w = 5$ and a threshold $\tau = 2$
\parencite{ouledmeriem2025hybrid}. Self-loops are added to every
node to ensure each item retains its own embedding during
propagation.

\subsection{Normalized Adjacency Matrix}

Let $\tilde{\mathbf{W}}$ denote the adjacency matrix of
$\mathcal{G}$ with self-loops, and let $\tilde{\mathbf{D}}$ be
the corresponding degree matrix with diagonal entries
$\tilde{D}_{ii} = \sum_j \tilde{W}_{ij}$. The
\textbf{symmetric normalized adjacency matrix} is computed as
\parencite{kipf2017gcn}:
\begin{equation}
  \hat{\mathbf{A}} = \tilde{\mathbf{D}}^{-1/2}\,
                     \tilde{\mathbf{W}}\,
                     \tilde{\mathbf{D}}^{-1/2}.
\end{equation}
This normalization prevents numerical instabilities during
propagation and ensures that items with many co-occurrence
connections do not dominate the aggregation. The resulting
$\hat{\mathbf{A}} \in \mathbb{R}^{|\mathcal{I}| \times |\mathcal{I}|}$
is sparse and symmetric.

\subsection{GCN Encoder}

Starting from the initial item embedding matrix
$\mathbf{E}^{(0)} = \mathbf{E} \in \mathbb{R}^{|\mathcal{I}| \times d}$,
we apply $L_{\text{GNN}}$ graph convolutional layers. At each
layer $\ell$, the item embeddings are updated by aggregating
the embeddings of neighboring items in the co-occurrence graph,
followed by a learnable linear transformation and a non-linear
activation:
\begin{equation}
  \mathbf{E}^{(\ell)} = \sigma\!\left(
    \hat{\mathbf{A}}\, \mathbf{E}^{(\ell-1)}\, \mathbf{W}^{(\ell)}
  \right), \quad \ell = 1, \ldots, L_{\text{GNN}},
\end{equation}
where $\mathbf{W}^{(\ell)} \in \mathbb{R}^{d \times d}$ is a
trainable weight matrix and $\sigma$ is the ReLU activation
function \parencite{ouledmeriem2025hybrid}. Each GCN layer
expands the receptive field of each item by one hop: after
$L_{\text{GNN}}$ layers, each item's embedding encodes
information from all items within $L_{\text{GNN}}$ hops in
the co-occurrence graph. In our experiments, we use
$L_{\text{GNN}} = 2$ layers, capturing up to two-hop
co-occurrence relationships.

The output of the final GCN layer,
$\mathbf{G} = \mathbf{E}^{(L_{\text{GNN}})} \in
\mathbb{R}^{|\mathcal{I}| \times d}$, provides
\textbf{graph-enhanced item embeddings}
$\mathbf{g}_i \in \mathbb{R}^d$ for each item $i$, encoding
the structural context of that item within the global
co-occurrence graph.

\subsection{Projection to BERT4Rec Embedding Space}

Since the GCN encoder and the BERT4Rec backbone may operate in
different embedding subspaces, a \textbf{linear projection layer}
is applied to map the graph-enhanced embeddings into the same
space as the BERT4Rec item embeddings
\parencite{ouledmeriem2025hybrid}:
\begin{equation}
  \mathbf{g}'_i = \mathbf{W}_g\, \mathbf{g}_i + \mathbf{b}_g,
\end{equation}
where $\mathbf{W}_g \in \mathbb{R}^{d \times d}$ and
$\mathbf{b}_g \in \mathbb{R}^d$ are learnable parameters. The
projected embeddings $\mathbf{g}'_i \in \mathbb{R}^d$ are then
ready to be fused with the BERT4Rec item embeddings through the
fusion mechanism described in Section~2.5. When the GCN and
BERT4Rec operate in the same embedding dimension $d$ --- as in
our experiments --- this projection still serves as a learned
alignment that calibrates the two embedding spaces before fusion.

\begin{table}[h]
\centering
\caption{GCN encoder hyperparameters}
\small
\begin{tabular}{lc}
\toprule
\textbf{Hyperparameter} & \textbf{Value} \\
\midrule
GCN layers $L_{\text{GNN}}$         & 2 \\
Co-occurrence window size $w$        & 5 \\
Min co-occurrence threshold $\tau$   & 2 \\
Embedding dimension $d$              & 64 \\
Activation function                  & ReLU \\
\bottomrule
\end{tabular}
\end{table}

% ============================================================
\section{The Four Hybrid Model Variants}
% ============================================================

All four variants proposed in this work share the same high-level
pipeline: the sequence path encodes $S_u$ via BERT4Rec producing
the sequence embedding $\mathbf{h}_u$; the auxiliary path produces
a complementary item-level signal; and the Fusion Layer blends
both into a unified item embedding $\mathbf{h}_i^{(u)}$ before the
final scoring step. The four variants differ exclusively in how
the auxiliary signal is constructed and how $\alpha(u)$ is
determined \parencite{ouledmeriem2025hybrid}.

% ------------------------------------------------------------
\subsection{Variant 1 — HybridFixed}
% ------------------------------------------------------------

\textbf{HybridFixed} is the simplest variant and serves as the
primary ablation reference within the hybrid family. It uses the
GCN graph embedding $\mathbf{g}'_i$ from Section~2.4 as the
auxiliary signal, and fuses it with the BERT4Rec item embedding
using a \textbf{single global fusion weight} $\alpha \in [0,1]$
shared identically across all users \parencite{ouledmeriem2025hybrid}.
For each item $i$, the fused embedding is:
\begin{equation}
  \mathbf{h}_i^{(u)} = \alpha\, \mathbf{e}_i
                      + (1 - \alpha)\, \mathbf{g}'_i,
\end{equation}
where $\alpha$ is a learnable scalar optimized end-to-end via
gradient descent alongside all other model parameters. This
allows $\alpha$ to converge to the globally optimal
GNN-Transformer trade-off for the full training distribution.

The purpose of HybridFixed is twofold: it verifies that combining
GNN and BERT4Rec embeddings is beneficial, and it serves as a
strong upper-bound reference that isolates the contribution of
the adaptive fusion mechanism. Despite its simplicity, HybridFixed
achieves the strongest \textit{aggregate} performance in our
experiments, improving HR@10 by $+8.1\%$, NDCG@10 by $+11.5\%$,
and MRR@10 by $+14.2\%$ over BERT4Rec on ML-1M
\parencite{ouledmeriem2025hybrid}.

% ------------------------------------------------------------
\subsection{Variant 2 — HybridDiscrete (Length-Adaptive)}
% ------------------------------------------------------------

\textbf{HybridDiscrete} is the core length-adaptive contribution
of this work. Rather than applying a single global $\alpha$ to
all users, it explicitly encodes the inductive bias that the
optimal fusion weight depends on the richness of each user's
interaction history \parencite{ouledmeriem2025hybrid}. Using the
same GCN auxiliary signal $\mathbf{g}'_i$ as HybridFixed, the
model assigns a \textbf{user-specific fusion weight} $\alpha(u)$
according to the three-bin partition defined in Section~2.2:
\begin{equation}
  \alpha(u) =
  \begin{cases}
    \alpha_{\text{short}} = 0.3 & \text{if } L(u) \leq 10 \\
    \alpha_{\text{mid}}   = 0.5 & \text{if } 10 < L(u) \leq 30 \\
    \alpha_{\text{long}}  = 0.7 & \text{if } L(u) > 30,
  \end{cases}
\end{equation}
and the fused item embedding is:
\begin{equation}
  \mathbf{h}_i^{(u)} = \alpha(u)\, \mathbf{e}_i
    + \bigl(1 - \alpha(u)\bigr)\, \mathbf{g}'_i.
\end{equation}

\textbf{Rationale.} A small $\alpha(u) = 0.3$ for short-history
users places $70\%$ of the weight on the GNN graph signal,
compensating for the insufficient sequential context available
to the Transformer. A balanced $\alpha(u) = 0.5$ for medium-history
users allows equal contribution from both signals. A large
$\alpha(u) = 0.7$ for long-history users preserves the full
expressiveness of the BERT4Rec sequence embedding. The three
$\alpha$ values are treated as hyperparameters tuned on the
validation set. This mechanism adds only three scalar parameters
and introduces negligible computational overhead
\parencite{ouledmeriem2025hybrid}.

\textbf{Why HybridDiscrete underperforms HybridFixed in aggregate.}
Although HybridDiscrete delivers the largest targeted gains for
short-history users, HybridFixed achieves stronger aggregate
metrics. This is explained by three factors: first, the
HybridFixed $\alpha$ is gradient-trained end-to-end and converges
to the globally optimal trade-off, while the discrete $\alpha$
values receive no gradient signal during training; second, the
discrete binning introduces hard boundary discontinuities between
users with $L(u) = 10$ and $L(u) = 11$, who receive sharply
different fusion weights despite nearly identical sequential
evidence; third, the Transformer in HybridFixed sees a single
globally consistent embedding space across all users, whereas
in HybridDiscrete it sees three distinct embedding spaces
depending on the user's bin \parencite{ouledmeriem2025hybrid}.
These limitations motivate a future direction of replacing the
discrete bins with a fully differentiable MLP that produces a
continuous, gradient-trained $\alpha(u)$.

% ------------------------------------------------------------
\subsection{Variant 3 — DCH-BERT4Rec}
% ------------------------------------------------------------

\textbf{DCH-BERT4Rec} (Dilated Convolutional Hybrid) replaces
the GCN graph encoder with a \textbf{dilated causal convolutional
network} \parencite{bai2018tcn} that encodes local sequential
transition patterns at multiple temporal scales. Rather than
building a global co-occurrence graph, this variant stacks
convolutional layers with kernel size $k$ and exponentially
increasing dilation rates:
\begin{equation}
  d_\ell = 2^\ell, \quad \ell = 1, 2, \ldots,
\end{equation}
producing a multi-scale receptive field that grows exponentially
with depth while maintaining causal (left-to-right) processing.
The convolutional output $\mathbf{c}_i \in \mathbb{R}^d$ captures
short-to-medium range co-occurrence patterns and is fused with
the BERT4Rec sequence embedding using the same fixed-weight
scheme as HybridFixed:
\begin{equation}
  \mathbf{h}_i^{(u)} = \alpha\, \mathbf{e}_i^{\text{BERT}}
                      + (1 - \alpha)\, \mathbf{c}_i.
\end{equation}

DCH-BERT4Rec tests whether \textit{local} convolutional transition
features provide complementary information to the \textit{global}
self-attention of BERT4Rec, without relying on a pre-built item
graph. Empirically, DCH-BERT4Rec achieves the strongest overall
results on ML-100K, improving HR@10 by $+25.8\%$ over BERT4Rec,
suggesting that dilated convolutions are particularly effective
on smaller and denser interaction graphs
\parencite{ouledmeriem2025hybrid}.

% ------------------------------------------------------------
\subsection{Variant 4 — TGT-BERT4Rec}
% ------------------------------------------------------------

\textbf{TGT-BERT4Rec} (Temporal Graph Transformer Hybrid)
augments the item co-occurrence graph with \textbf{temporal
information} by encoding interaction timestamps as continuous
time embeddings \parencite{ouledmeriem2025hybrid}. A temporal
attention bias is added to the standard scaled dot-product
attention, allowing the model to modulate attention scores
based on the time gap between interactions:
\begin{equation}
  \text{Attention}(\mathbf{Q}, \mathbf{K}, \mathbf{V}, \mathbf{T}) =
    \text{softmax}\!\left(
      \frac{\mathbf{Q}\mathbf{K}^\top + \text{TimeBias}(\mathbf{T})}
           {\sqrt{d_k}}
    \right) \mathbf{V},
\end{equation}
where $\text{TimeBias}(\mathbf{T})$ is a learned function of the
time-gap embedding $\mathbf{T}$ between interactions
\parencite{ouledmeriem2025hybrid}. This enriches the item
representations with temporal context beyond simple positional
encoding.

A \textbf{learnable gated fusion} combines the TGT encoder output
$\mathbf{h}_i^{\text{TGT}}$ with the BERT4Rec sequence
representation:
\begin{equation}
  \alpha = \sigma(g), \qquad
  \mathbf{h}_i^{(u)} = \alpha\, \mathbf{h}_i^{\text{BERT}}
                      + (1 - \alpha)\, \mathbf{h}_i^{\text{TGT}},
\end{equation}
where $g$ is a learnable scalar gate and $\sigma$ is the sigmoid
function. This gating mechanism allows the model to learn the
optimal balance between the global self-attention representation
and the temporal graph encoding. TGT-BERT4Rec is particularly
well-suited to short-history users, whose temporal interaction
patterns encode strong preference signals even when interaction
counts are low. Empirically, it achieves the highest HR@10 for
short-history users on both ML-100K and ML-1M, with HR@10 of
$0.2778$ and $0.2209$ respectively
\parencite{ouledmeriem2025hybrid}.

% ------------------------------------------------------------
\subsection{Summary of the Four Variants}
% ------------------------------------------------------------

\begin{table}[h]
\centering
\caption{Summary of the four hybrid model variants}
\renewcommand{\arraystretch}{1.3}
\footnotesize
\begin{tabular}{p{3.2cm} p{3.2cm} p{3.2cm} p{2.8cm}}
\toprule
\textbf{Variant} & \textbf{Auxiliary signal} &
\textbf{Fusion type} & \textbf{Best on} \\
\midrule
HybridFixed     & GCN co-occurrence graph  & Global learnable $\alpha$ & ML-1M overall \\
HybridDiscrete  & GCN co-occurrence graph  & Per-user discrete $\alpha(u)$ & Short-history users \\
DCH-BERT4Rec    & Dilated conv encoder     & Global learnable $\alpha$ & ML-100K overall \\
TGT-BERT4Rec    & Temporal graph encoder   & Gated sigmoid $\alpha$   & Short-history users \\
\bottomrule
\end{tabular}
\end{table}

% ============================================================
\section{Training Objective}
% ============================================================

\subsection{Binary Cross-Entropy Loss}

All four hybrid variants are trained with a
\textbf{point-wise Binary Cross-Entropy (BCE) loss} with sampled
negative items \parencite{ouledmeriem2025hybrid}. For each
training instance $(u, t)$, let $i^+$ denote the ground-truth
positive item at time step $t+1$, and let $\mathcal{N}_{u,t}$
denote a set of randomly sampled negative items --- items that
user $u$ has not interacted with. The relevance score for any
candidate item $j$ is computed as:
\begin{equation}
  s_{uj} = \mathbf{h}_u^\top\, \mathbf{h}_j^{(u)},
\end{equation}
where $\mathbf{h}_u$ is the sequence embedding from BERT4Rec
and $\mathbf{h}_j^{(u)}$ is the fused item embedding from the
fusion layer. The training loss is:
\begin{equation}
  \mathcal{L} = -\sum_{(u,t)}
    \left[
      \log \sigma\!\left(s_{u i^+}\right)
      + \sum_{j \in \mathcal{N}_{u,t}}
        \log\!\left(1 - \sigma\!\left(s_{uj}\right)\right)
    \right]
    + \lambda \|\Theta\|_2^2,
\end{equation}
where $\sigma(\cdot)$ is the sigmoid function, $\lambda$ is the
$L_2$ regularization coefficient, and $\Theta$ denotes all
trainable parameters \parencite{ouledmeriem2025hybrid}. The
regularization term penalizes large parameter norms and reduces
overfitting, which is particularly important for sparse users
with few training instances.

\subsection{Negative Sampling Strategy}

For each positive training instance $(u, i^+)$, we sample
\textbf{100 negative items} uniformly at random from
$\mathcal{I} \setminus S_u$ --- the set of items not present
in user $u$'s interaction history \parencite{ouledmeriem2025hybrid}.
This large number of negatives per positive provides a rich
contrastive training signal and ensures that the model learns
to distinguish the ground-truth item from a diverse set of
distractors at every training step. BCE with sampled negatives
is preferred over full softmax over the item vocabulary for
computational efficiency, as the latter requires computing
scores for all $|\mathcal{I}|$ items at every training step.

\subsection{Joint End-to-End Optimization}

All model components are optimized \textbf{jointly end-to-end}
using a single loss function \parencite{ouledmeriem2025hybrid}.
This includes the BERT4Rec Transformer parameters, the GCN
encoder weight matrices $\{\mathbf{W}^{(\ell)}\}$, the projection
layer $(\mathbf{W}_g, \mathbf{b}_g)$, and the fusion weights.
Joint optimization ensures that the GCN and BERT4Rec components
co-adapt during training: the GCN learns to produce graph
embeddings that are maximally complementary to the Transformer's
sequential representations, rather than being frozen or
pre-trained independently.

The sole exception is the HybridDiscrete variant, where the three
scalar values $(\alpha_{\text{short}}, \alpha_{\text{mid}},
\alpha_{\text{long}})$ are treated as hyperparameters tuned on
the validation set rather than gradient-trained parameters.
As discussed in Section~2.5, this means these values receive no
gradient signal during training, which partially explains
HybridDiscrete's lower aggregate performance compared to
HybridFixed \parencite{ouledmeriem2025hybrid}.

\subsection{Optimizer and Training Protocol}

All models are optimized using the \textbf{Adam optimizer}
\parencite{kingma2015adam} with the following settings:
\begin{itemize}
  \item Learning rate: $\eta = 0.001$
  \item Batch size: $256$
  \item Maximum epochs: $200$
  \item Early stopping patience: $20$ epochs on validation NDCG@10
\end{itemize}

Training is stopped when the validation NDCG@10 does not improve
for 20 consecutive epochs, and the model checkpoint with the best
validation NDCG@10 is used for final test evaluation. This
protocol ensures fair comparison across all models and prevents
overfitting to the training set \parencite{ouledmeriem2025hybrid}.

% ============================================================
\section{Hyperparameters and Complexity Analysis}
% ============================================================

\subsection{Hyperparameter Settings}

All models are implemented and evaluated under identical
conditions for fair comparison. Table~\ref{tab:hyperparams}
summarizes the full set of hyperparameters used across all
experiments \parencite{ouledmeriem2025hybrid}.

\begin{table}[h]
\centering
\caption{Hyperparameter settings for all models}
\label{tab:hyperparams}
\renewcommand{\arraystretch}{1.3}
\small
\begin{tabular}{lc}
\toprule
\textbf{Parameter} & \textbf{Value} \\
\midrule
\multicolumn{2}{l}{\textit{General}} \\
Embedding dimension $d$              & 64 \\
Attention heads                      & 2 \\
Transformer blocks $B$               & 2 \\
Dropout rate                         & 0.2 \\
Max sequence length (hybrid models)  & 200 \\
Max sequence length (baselines)      & 50 \\
\midrule
\multicolumn{2}{l}{\textit{Training}} \\
Optimizer                            & Adam \\
Learning rate $\eta$                 & 0.001 \\
Batch size                           & 256 \\
Max epochs                           & 200 \\
Early stopping patience              & 20 (val NDCG@10) \\
Negative samples per positive        & 100 \\
\midrule
\multicolumn{2}{l}{\textit{GCN Encoder}} \\
GCN layers $L_{\text{GNN}}$          & 2 \\
Co-occurrence window size $w$        & 5 \\
Min co-occurrence threshold $\tau$   & 2 \\
\midrule
\multicolumn{2}{l}{\textit{HybridDiscrete Fusion}} \\
Short-history threshold $L^*_{\text{short}}$ & 10 \\
Long-history threshold $L^*_{\text{long}}$   & 30 \\
$\alpha_{\text{short}}$              & 0.3 \\
$\alpha_{\text{mid}}$                & 0.5 \\
$\alpha_{\text{long}}$               & 0.7 \\
\bottomrule
\end{tabular}
\end{table}

\noindent One notable asymmetry in Table~\ref{tab:hyperparams}
is the difference in maximum sequence length between hybrid models
(200) and pure baselines (50). This difference exists because
hybrid models leverage the GCN graph signal to compensate for
sparse users, making longer sequences informative without
introducing noise. However, this is acknowledged as a
methodological limitation: future work should control for uniform
truncation across all model families to ensure strictly fair
comparison \parencite{ouledmeriem2025hybrid}.

\subsection{Complexity Analysis}

\textbf{Time complexity.} The dominant computational cost of the
proposed hybrid models comes from two components. The GCN encoder
performs $L_{\text{GNN}}$ sparse matrix multiplications over the
co-occurrence graph, with time complexity
$\mathcal{O}(|\mathcal{E}| \cdot d)$ per layer, where
$|\mathcal{E}|$ is the number of edges in $\mathcal{G}$ and $d$
is the embedding dimension. The BERT4Rec Transformer applies
self-attention over sequences of length $L$, with time complexity
$\mathcal{O}(L^2 \cdot d)$ per attention block. The fusion
operation adds only $\mathcal{O}(|\mathcal{I}| \cdot d)$
element-wise operations, which is negligible compared to the
above. The total time complexity per training step is therefore:
\begin{equation}
  \mathcal{O}\!\left(
    L_{\text{GNN}} \cdot |\mathcal{E}| \cdot d
    + B \cdot L^2 \cdot d
  \right),
\end{equation}
which is dominated by the Transformer's quadratic attention when
$L$ is large \parencite{ouledmeriem2025hybrid}.

\textbf{Parameter complexity.} A key strength of the proposed
architecture is its \textbf{parameter efficiency}. All four
hybrid variants add only approximately $0.1$M parameters over
the BERT4Rec backbone, corresponding to a $+5.9\%$ increase
from $1.7$M to $1.8$M total parameters
\parencite{ouledmeriem2025hybrid}. The additional parameters
correspond exclusively to the GCN weight matrices
$\{\mathbf{W}^{(\ell)}\}$ and the projection layer
$(\mathbf{W}_g, \mathbf{b}_g)$. This confirms that the
performance gains reported in Chapter~3 are driven by
\textbf{architectural complementarity} rather than by an
increased parameter budget. Table~\ref{tab:params} summarizes
the parameter counts for all models.

\begin{table}[h]
\centering
\caption{Trainable parameter counts for all models on ML-1M}
\label{tab:params}
\renewcommand{\arraystretch}{1.3}
\small
\begin{tabular}{lcc}
\toprule
\textbf{Model} & \textbf{\# Parameters} & \textbf{vs BERT4Rec} \\
\midrule
GRU4Rec                  & 1.2M & --- \\
LightGCN                 & 0.8M & --- \\
Caser                    & 1.1M & --- \\
SASRec                   & 1.7M & --- \\
BERT4Rec                 & 1.7M & baseline \\
\midrule
DCH-BERT4Rec\textsuperscript{†}    & 1.8M & +5.9\% \\
HybridDiscrete\textsuperscript{†}  & 1.8M & +5.9\% \\
TGT-BERT4Rec\textsuperscript{†}    & 1.8M & +5.9\% \\
HybridFixed\textsuperscript{†}     & 1.8M & +5.9\% \\
\bottomrule
\multicolumn{3}{l}{\textsuperscript{†}Proposed hybrid model}
\end{tabular}
\end{table}

% ============================================================
\section{Chapter Summary}
% ============================================================

This chapter presented the methodology of the first contribution
of this memoir: the \textbf{Length-Adaptive Hybrid GNN-Transformer}
for sequential recommendation \parencite{ouledmeriem2025hybrid}.
The central motivation was the observation that existing hybrid
approaches apply a fixed fusion weight uniformly to all users,
ignoring a fundamental reality of interaction data: the relative
reliability of graph signals versus sequential signals varies
dramatically with the length of each user's interaction history.

The proposed architecture addresses this limitation through a
shared pipeline combining two complementary encoders. The
\textbf{BERT4Rec backbone} captures fine-grained temporal
dynamics within each user's interaction sequence through
bidirectional self-attention. The \textbf{GCN encoder} captures
global structural relationships between items by propagating
embeddings over an item co-occurrence graph built from all
training sequences. These two signals are blended through a
\textbf{fusion layer} that produces a unified item embedding
$\mathbf{h}_i^{(u)} = \alpha(u)\,\mathbf{e}_i +
(1-\alpha(u))\,\mathbf{g}'_i$, where $\alpha(u) \in [0,1]$
controls the trade-off between the two sources of information.

Four variants were introduced, each differing in how the
auxiliary signal is constructed and how $\alpha(u)$ is
determined:

\begin{itemize}
  \item \textbf{HybridFixed} applies a single global learnable
        $\alpha$ to all users, achieving the strongest aggregate
        performance through end-to-end gradient optimization.

  \item \textbf{HybridDiscrete} applies a user-specific
        $\alpha(u)$ determined by three discrete bins based on
        interaction history length, delivering the largest
        targeted gains for short-history and sparse users.

  \item \textbf{DCH-BERT4Rec} replaces the GCN with a dilated
        causal convolutional encoder, capturing local sequential
        transition patterns at multiple temporal scales.

  \item \textbf{TGT-BERT4Rec} augments the co-occurrence graph
        with temporal interaction timestamps and uses a gated
        fusion mechanism, achieving the strongest performance
        for short-history users on both datasets.
\end{itemize}

All four variants share the same training objective --- Binary
Cross-Entropy loss with 100 sampled negatives per positive,
optimized end-to-end with Adam --- and add only $+0.1$M
parameters ($+5.9\%$) over the BERT4Rec backbone, confirming
that their performance gains stem from architectural
complementarity rather than increased model capacity
\parencite{ouledmeriem2025hybrid}.

The experimental evaluation of all four variants against the
five baselines (GRU4Rec, Caser, SASRec, BERT4Rec, LightGCN)
on MovieLens-100K and MovieLens-1M is presented in the
following chapter.



% ============================================================
\chapter{Experimental Evaluation}
\addcontentsline{toc}{chapter}{Chapter 3 — Experimental Evaluation}
% ============================================================

\section{Datasets}

All experiments are conducted on two standard sequential recommendation benchmarks from the MovieLens collection \parencite{harper2015movielens}: \textbf{MovieLens-100K} (ML-100K) and \textbf{MovieLens-1M} (ML-1M). These datasets are widely used in the sequential recommendation literature \parencite{kang2018self, sun2019bert4rec, he2020lightgcn} and provide a rigorous basis for comparison with existing work.

\textbf{MovieLens-100K} contains 100,000 ratings from 943 users on 1,682 movies. After filtering users with fewer than 5 interactions and converting ratings to implicit feedback, the dataset yields \textbf{938 test users}.

\textbf{MovieLens-1M} contains approximately 1 million ratings from 6,040 users on 3,706 movies. After applying the same preprocessing, the dataset yields \textbf{6,034 test users}.

\subsection{Preprocessing and Evaluation Protocol}

For both datasets, we adopt the \textbf{leave-one-out} evaluation protocol \parencite{kang2018self}: the last item in each user's chronologically ordered interaction sequence is held out for testing, the second-to-last item for validation, and all remaining items constitute the training set. Interaction sequences are truncated to a maximum length of 200 for hybrid models and 50 for pure sequential baselines, following the implementation details in \parencite{ouledmeriem2025hybrid}.

During evaluation, each ground-truth test item is ranked against \textbf{all items} in the full catalog, without negative sampling. This is a stricter evaluation protocol than the sampled-negative approach and provides a more reliable measure of true ranking performance.

\section{Implementation Details and Training Protocol}

All models are implemented under identical conditions for fair comparison. Table~\ref{tab:hyperparams_ch3} reproduces the complete hyperparameter settings used across all experiments \parencite{ouledmeriem2025hybrid}.

\begin{table}[h]
\centering
\caption{Hyperparameter settings for all models}
\label{tab:hyperparams_ch3}
\renewcommand{\arraystretch}{1.3}
\adjustbox{max width=\textwidth}{%
\begin{tabular}{lccccccccc}
\toprule
\textbf{Model} & \textbf{d} & \textbf{heads} & \textbf{blocks} & \textbf{layers} & \textbf{dropout} & \textbf{max\_len} & \textbf{lr} & \textbf{batch} & \textbf{patience} \\
\midrule
BERT4Rec             & 64 & 2 & 2 & 1 & 0.2 & 50  & 0.001 & 256 & 20 \\
BERT-Hybrid-Discrete & 64 & 2 & 2 & — & 0.2 & 200 & 0.001 & 256 & 20 \\
BERT-Hybrid-Fixed    & 64 & 2 & 2 & — & 0.2 & 200 & 0.001 & 256 & 20 \\
Caser                & 64 & 2 & 2 & 1 & 0.2 & 50  & 0.001 & 256 & 20 \\
GRU4Rec              & 64 & 2 & 2 & 1 & 0.2 & 50  & 0.001 & 256 & 20 \\
LightGCN             & 64 & 2 & 2 & 1 & 0.2 & 50  & 0.001 & 256 & 20 \\
SASRec               & 64 & 2 & 2 & 1 & 0.2 & 50  & 0.001 & 256 & 20 \\
TCN-BERT4Rec         & 64 & 2 & 2 & — & 0.2 & 200 & 0.001 & 256 & 20 \\
TGT-BERT4Rec         & 64 & 2 & 2 & — & 0.2 & 200 & 0.001 & 256 & 20 \\
\bottomrule
\end{tabular}%
}
\end{table}

All models are optimized using the Adam optimizer with learning rate \(\eta = 0.001\), batch size 256, and early stopping with patience 20 on validation NDCG@10. The best model checkpoint is selected based on the highest validation NDCG@10 across up to 200 epochs \parencite{ouledmeriem2025hybrid}.

Table~\ref{tab:best_epoch} reports the best epoch and best validation NDCG@10 for each model. BERT4Rec converges at epoch 50 with a validation NDCG@10 of 0.0705, while TGT-BERT4Rec requires the most training, reaching its best checkpoint at epoch 148. GRU4Rec converges fastest at epoch 30 with the lowest validation NDCG@10 of 0.0335, consistent with its weakest test performance. LightGCN converges at epoch 1, reflecting the absence of any sequential modeling capacity.

\begin{table}[h]
\centering
\caption{Best epoch and validation NDCG@10 per model on ML-1M}
\label{tab:best_epoch}
\renewcommand{\arraystretch}{1.3}
\small
\begin{tabular}{lcc}
\toprule
\textbf{Model} & \textbf{Best Epoch} & \textbf{Best Val NDCG@10} \\
\midrule
BERT4Rec             & 50  & 0.0705 \\
BERT-Hybrid-Discrete & 64  & — \\
BERT-Hybrid-Fixed    & 88  & — \\
Caser                & 63  & 0.0657 \\
GRU4Rec              & 30  & 0.0335 \\
LightGCN             & 1   & 0.0107 \\
SASRec               & 113 & 0.0531 \\
TCN-BERT4Rec         & 84  & — \\
TGT-BERT4Rec         & 148 & — \\
\bottomrule
\multicolumn{3}{l}{— denotes full-catalog evaluation (no sampled validation)} \\
\end{tabular}
\end{table}

\subsection{Training Dynamics}

Figure~\ref{fig:train_loss} shows the training loss per epoch for all models on ML-1M. All models except LightGCN exhibit a rapid decrease in the first 25 epochs followed by a gradual plateau, confirming stable convergence. LightGCN exhibits a flat training loss from the very first epoch, reflecting that its purely graph-based collaborative filtering objective converges immediately without sequential modeling. GRU4Rec shows a slower initial descent and a higher plateau than all BERT-based models, consistent with the known difficulty of training recurrent architectures on long-tail interaction distributions. TGT-BERT4Rec achieves the lowest final training loss, suggesting that temporal graph signals provide the richest training objective.

\begin{figure}[h]
\centering
\includegraphics[width=\textwidth]{picture/figures/training_loss_1m.png}
\caption{Training loss per epoch for all models on ML-1M. LightGCN converges at epoch 1; TGT-BERT4Rec achieves the lowest final loss.}
\label{fig:train_loss}
\end{figure}

Figure~\ref{fig:val_ndcg} shows the validation NDCG@10 per epoch. The BERT-based hybrid models cluster near the top of the chart, all achieving validation NDCG@10 above 0.063. SASRec lags behind all BERT-based models despite converging later (epoch 113), confirming the advantage of bidirectional attention. GRU4Rec plateaus early at a low validation NDCG@10 of 0.033, limited by the vanishing gradient problem on long sequences \parencite{hidasi2016gru4rec}.

\begin{figure}[h]
\centering
\includegraphics[width=\textwidth]{picture/figures/val_ndgc_10.png}
\caption{Validation NDCG@10 per epoch for all models on ML-1M. Dots mark the best checkpoint used for final evaluation. TGT-BERT4Rec achieves the highest peak validation NDCG@10 (0.0745) at epoch 148.}
\label{fig:val_ndcg}
\end{figure}

\section{Baselines}

We compare the four proposed hybrid variants against five established baselines that represent the major paradigms in sequential recommendation \parencite{ouledmeriem2025hybrid}:

\begin{itemize}
  \item \textbf{GRU4Rec} \parencite{hidasi2016gru4rec}: The first recurrent model for session-based recommendation, using Gated Recurrent Units to encode temporal interaction sequences.
  \item \textbf{Caser} \parencite{tang2018caser}: A convolutional sequence model that captures point-level and union-level sequential patterns via horizontal and vertical filters over the most recent items.
  \item \textbf{SASRec} \parencite{kang2018self}: A unidirectional Transformer with causal self-attention, trained with binary cross-entropy loss.
  \item \textbf{BERT4Rec} \parencite{sun2019bert4rec}: A bidirectional Transformer trained with a Cloze-style masked item prediction objective. This is the primary baseline and the shared backbone of all hybrid variants.
  \item \textbf{LightGCN} \parencite{he2020lightgcn}: A simplified graph convolutional network that performs pure neighborhood aggregation on the user-item bipartite graph. It is included as a reference point to isolate the contribution of pure collaborative filtering signals.
\end{itemize}

\section{Main Results on MovieLens-1M}

Table~\ref{tab:ml1m_full} reports the test performance of all nine models on ML-1M across HR, NDCG, and MRR at cutoffs \(K \in \{5, 10, 20\}\). Bold values indicate the best result per column; underlined values indicate the second-best.

\begin{table}[h]
\centering
\caption{Test performance on MovieLens-1M. Best results are bold, second-best underlined.}
\label{tab:ml1m_full}
\renewcommand{\arraystretch}{1.3}
\adjustbox{max width=\textwidth}{%
\begin{tabular}{lrrrrrrrrr}
\toprule
\textbf{Model} &
\textbf{HR@5} & \textbf{NDCG@5} & \textbf{MRR@5} &
\textbf{HR@10} & \textbf{NDCG@10} & \textbf{MRR@10} &
\textbf{HR@20} & \textbf{NDCG@20} & \textbf{MRR@20} \\
\midrule
BERT4Rec
  & 0.0757 & 0.0453 & 0.0354
  & 0.1374 & 0.0652 & 0.0437
  & 0.2274 & 0.0879 & 0.0498 \\
BERT-Hybrid-Discrete
  & 0.0832 & 0.0492 & 0.0381
  & 0.1427 & 0.0682 & 0.0459
  & 0.2317 & 0.0905 & 0.0519 \\
BERT-Hybrid-Fixed
  & \textbf{0.0852} & \textbf{0.0524} & \textbf{0.0417}
  & \textbf{0.1485} & \textbf{0.0727} & \textbf{0.0499}
  & \textbf{0.2511} & \textbf{0.0984} & \textbf{0.0569} \\
Caser
  & 0.0684 & 0.0397 & 0.0303
  & 0.1260 & 0.0581 & 0.0379
  & 0.2135 & 0.0801 & 0.0438 \\
GRU4Rec
  & 0.0365 & 0.0214 & 0.0165
  & 0.0708 & 0.0324 & 0.0210
  & 0.1173 & 0.0440 & 0.0241 \\
LightGCN
  & 0.0121 & 0.0065 & 0.0047
  & 0.0247 & 0.0106 & 0.0064
  & 0.0467 & 0.0161 & 0.0079 \\
SASRec
  & 0.0515 & 0.0298 & 0.0228
  & 0.0979 & 0.0446 & 0.0287
  & 0.1767 & 0.0643 & 0.0340 \\
TCN-BERT4Rec
  & 0.0754 & 0.0449 & 0.0350
  & 0.1385 & 0.0652 & 0.0433
  & 0.2325 & 0.0888 & 0.0498 \\
TGT-BERT4Rec
  & \underline{0.0820} & \underline{0.0477} & \underline{0.0365}
  & \underline{0.1480} & \underline{0.0689} & \underline{0.0452}
  & \underline{0.2444} & \underline{0.0930} & \underline{0.0517} \\
\bottomrule
\end{tabular}%
}
\end{table}

\subsection{Overall Performance Analysis}

The results in Table~\ref{tab:ml1m_full} confirm three clear findings. First, \textbf{BERT-Hybrid-Fixed achieves the strongest aggregate performance} across all nine metrics, improving over BERT4Rec by \(+8.1\%\) on HR@10, \(+11.5\%\) on NDCG@10, and \(+14.2\%\) on MRR@10 \parencite{ouledmeriem2025hybrid}. Second, all four hybrid variants outperform all five baselines on at least one metric, confirming that augmenting BERT4Rec with a complementary structural or temporal signal yields consistent gains. Third, LightGCN performs worst by a significant margin, confirming that static graph-based collaborative filtering without temporal modeling is insufficient for next-item prediction tasks.

The heatmap in Figure~\ref{fig:heatmap_1m} provides a visual overview of all metrics simultaneously. The consistently deeper red shading of BERT-Hybrid-Fixed and TGT-BERT4Rec across all columns confirms their superiority, while the uniformly pale coloring of LightGCN and GRU4Rec quantifies the gap between graph-only and sequence-aware architectures.

\begin{figure}[h]
\centering
\includegraphics[width=\textwidth]{picture/figures/heat_map_1m.png}
\caption{All test metrics heatmap on ML-1M. Darker red indicates higher metric values. BERT-Hybrid-Fixed and TGT-BERT4Rec consistently show the deepest shading.}
\label{fig:heatmap_1m}
\end{figure}

\subsection{Relative Improvement over BERT4Rec}

Table~\ref{tab:improvements} reports the relative improvement (\%) of each model over BERT4Rec across all metrics at \(K \in \{5, 10, 20\}\). BERT-Hybrid-Fixed achieves the largest gains at every cutoff and every metric: \(+12.4\%\) HR@5, \(+15.7\%\) NDCG@5, \(+17.8\%\) MRR@5, \(+8.1\%\) HR@10, \(+11.5\%\) NDCG@10, and \(+14.3\%\) MRR@10. TGT-BERT4Rec is the strongest runner-up with \(+7.7\%\) HR@10. In contrast, GRU4Rec, LightGCN, and SASRec all show large negative improvements relative to BERT4Rec, confirming the superiority of bidirectional attention over recurrent and unidirectional architectures.

\begin{table}[h]
\centering
\caption{Relative improvement (\%) compared to BERT4Rec on ML-1M.}
\label{tab:improvements}
\renewcommand{\arraystretch}{1.3}
\adjustbox{max width=\textwidth}{%
\begin{tabular}{lrrrrrrrrr}
\toprule
\textbf{Model} &
$\Delta$ \textbf{HR@5} & $\Delta$ \textbf{NDCG@5} & $\Delta$ \textbf{MRR@5} &
$\Delta$ \textbf{HR@10} & $\Delta$ \textbf{NDCG@10} & $\Delta$ \textbf{MRR@10} &
$\Delta$ \textbf{HR@20} & $\Delta$ \textbf{NDCG@20} & $\Delta$ \textbf{MRR@20} \\
\midrule
BERT4Rec             & 0.00 & 0.00 & 0.00 & 0.00 & 0.00 & 0.00 & 0.00 & 0.00 & 0.00 \\
BERT-Hybrid-Discrete & +9.85 & +8.52 & +7.50 & +3.86 & +4.57 & +5.02 & +1.90 & +2.98 & +4.08 \\
BERT-Hybrid-Fixed    & \textbf{+12.47} & \textbf{+15.73} & \textbf{+17.77} & \textbf{+8.08} & \textbf{+11.37} & \textbf{+14.38} & \textbf{+10.42} & \textbf{+11.97} & \textbf{+14.22} \\
Caser                & $-9.63$ & $-12.47$ & $-14.20$ & $-8.32$ & $-10.90$ & $-13.26$ & $-6.12$ & $-8.84$ & $-12.01$ \\
GRU4Rec              & $-51.86$ & $-52.85$ & $-53.49$ & $-48.49$ & $-50.30$ & $-51.84$ & $-48.40$ & $-49.92$ & $-51.61$ \\
LightGCN             & $-84.03$ & $-85.58$ & $-86.64$ & $-82.03$ & $-83.83$ & $-85.43$ & $-79.45$ & $-81.71$ & $-84.24$ \\
SASRec               & $-31.95$ & $-34.14$ & $-35.70$ & $-28.71$ & $-31.60$ & $-34.17$ & $-22.30$ & $-26.86$ & $-31.73$ \\
TCN-BERT4Rec         & $-0.44$ & $-0.86$ & $-1.11$ & $+0.84$ & $-0.07$ & $-0.76$ & $+2.26$ & $+1.07$ & $-0.15$ \\
TGT-BERT4Rec         & $+8.32$ & $+5.25$ & $+3.06$ & $+7.72$ & $+5.58$ & $+3.51$ & $+7.51$ & $+5.83$ & $+3.75$ \\
\bottomrule
\end{tabular}%
}
\end{table}

\section{Per-History-Length Analysis}

A core motivation of the length-adaptive design is that the relative utility of graph versus sequential signals varies with the user's interaction history length. To validate this, we segment all 6,034 test users on ML-1M into three groups: \textbf{SHORT} (\(L(u) \leq 10\), \(n = 162\) users), \textbf{MEDIUM} (\(10 < L(u) \leq 30\), \(n = 2{,}579\) users), and \textbf{OVERALL} (\(n = 6{,}034\) users) \parencite{ouledmeriem2025hybrid}.

\subsection{Short-History Users}

For \textbf{short-history users} (Table~\ref{tab:bucket_short}), TGT-BERT4Rec achieves the best HR@10 (0.2778) and NDCG@10 (0.1213), outperforming BERT4Rec by \(+18.4\%\) and \(+11.8\%\) respectively. This confirms that temporal graph signals are particularly informative when sequential evidence is scarce: even with fewer than 10 interactions, the time gaps between interactions encode strong preference signals. BERT-Hybrid-Discrete achieves the same HR@10 as BERT-Hybrid-Fixed (0.2469) for short-history users, directly validating the length-adaptive design: by assigning \(\alpha = 0.3\) to short-history users, the model compensates for insufficient sequential context by leaning on the GNN graph signal.

Figures~\ref{fig:hr10_buckets}, \ref{fig:ndcg10_buckets}, and \ref{fig:mrr10_buckets} visualize HR@10, NDCG@10, and MRR@10 across the three groups for all models. The consistent pattern across all three metrics confirms that TGT-BERT4Rec dominates for short-history users and BERT-Hybrid-Fixed dominates for medium-history users.

\begin{figure}[h]
\centering
\includegraphics[width=\textwidth]{picture/figures/hr10_buckets.png}
\caption{HR@10 across Short, Medium, and Overall user groups on ML-1M. TGT-BERT4Rec achieves the highest short-history HR@10 (0.2778); BERT-Hybrid-Fixed leads for medium and overall.}
\label{fig:hr10_buckets}
\end{figure}

\begin{figure}[h]
\centering
\includegraphics[width=\textwidth]{picture/figures/ndcg10_buckets.png}
\caption{NDCG@10 across Short, Medium, and Overall user groups on ML-1M. TGT-BERT4Rec achieves NDCG@10 of 0.1213 for short-history users, a \(+11.8\%\) improvement over BERT4Rec.}
\label{fig:ndcg10_buckets}
\end{figure}

\begin{figure}[h]
\centering
\includegraphics[width=\textwidth]{picture/figures/mrr10_buckets.png}
\caption{MRR@10 across Short, Medium, and Overall user groups on ML-1M. All hybrid models outperform BERT4Rec across all three user groups.}
\label{fig:mrr10_buckets}
\end{figure}

\subsection{Medium-History Users}

For \textbf{medium-history users} (\(2{,}579\) users with \(10 < L(u) \leq 30\)), BERT-Hybrid-Fixed achieves the top HR@10 (0.1927), representing a \(+43.1\%\) improvement over BERT4Rec (0.1347). TGT-BERT4Rec ranks second with HR@10 of 0.1912. BERT-Hybrid-Discrete also substantially improves over BERT4Rec with HR@10 of 0.1834 (\(+36.2\%\)), demonstrating that all hybrid variants generalize well to medium-length sequences. Figures~\ref{fig:hr20_buckets}, \ref{fig:ndcg20_buckets}, and \ref{fig:mrr20_buckets} extend this analysis to the @20 cutoff, showing identical ranking patterns.

\begin{figure}[h]
\centering
\includegraphics[width=\textwidth]{picture/figures/hr20_buckets.png}
\caption{HR@20 across Short, Medium, and Overall user groups on ML-1M. TGT-BERT4Rec leads for short-history users (0.3642), BERT-Hybrid-Fixed leads for medium (0.3110).}
\label{fig:hr20_buckets}
\end{figure}

\begin{figure}[h]
\centering
\includegraphics[width=\textwidth]{picture/figures/ndcg20_buckets.png}
\caption{NDCG@20 across Short, Medium, and Overall user groups on ML-1M.}
\label{fig:ndcg20_buckets}
\end{figure}

\begin{figure}[h]
\centering
\includegraphics[width=\textwidth]{picture/figures/mrr20_buckets.png}
\caption{MRR@20 across Short, Medium, and Overall user groups on ML-1M.}
\label{fig:mrr20_buckets}
\end{figure}

\begin{table}[h]
\centering
\caption{Per-bucket HR@10 / NDCG@10 / MRR@10 on ML-1M.}
\label{tab:bucket_short}
\renewcommand{\arraystretch}{1.3}
\adjustbox{max width=\textwidth}{%
\begin{tabular}{lrrr|rrr|rrr}
\toprule
& \multicolumn{3}{c|}{\textbf{SHORT} (\(n=162\))} & \multicolumn{3}{c|}{\textbf{MEDIUM} (\(n=2{,}579\))} & \multicolumn{3}{c}{\textbf{OVERALL} (\(n=6{,}034\))} \\
\textbf{Model}
  & HR@10 & NDCG@10 & MRR@10 & HR@10 & NDCG@10 & MRR@10 & HR@10 & NDCG@10 & MRR@10 \\
\midrule
BERT4Rec
  & 0.2346 & 0.1085 & 0.0707 & 0.1347 & 0.0640 & 0.0429 & 0.1374 & 0.0652 & 0.0437 \\
BERT-Hybrid-Discrete
  & 0.2469 & 0.1132 & 0.0732 & 0.1834 & 0.0884 & 0.0598 & 0.1427 & 0.0682 & 0.0459 \\
BERT-Hybrid-Fixed
  & 0.2469 & 0.1102 & 0.0700 & \textbf{0.1927} & \textbf{0.0963} & \textbf{0.0672} & \textbf{0.1485} & \textbf{0.0727} & \textbf{0.0499} \\
Caser
  & 0.2284 & 0.0964 & 0.0580 & 0.1231 & 0.0571 & 0.0373 & 0.1260 & 0.0581 & 0.0379 \\
GRU4Rec
  & 0.0494 & 0.0229 & 0.0149 & 0.0714 & 0.0327 & 0.0212 & 0.0708 & 0.0324 & 0.0210 \\
LightGCN
  & 0.0309 & 0.0125 & 0.0071 & 0.0245 & 0.0105 & 0.0063 & 0.0247 & 0.0106 & 0.0064 \\
SASRec
  & 0.1420 & 0.0776 & 0.0583 & 0.0967 & 0.0437 & 0.0279 & 0.0979 & 0.0446 & 0.0287 \\
TCN-BERT4Rec
  & 0.2222 & 0.1066 & 0.0710 & 0.1838 & 0.0867 & 0.0578 & 0.1385 & 0.0652 & 0.0433 \\
TGT-BERT4Rec
  & \textbf{0.2778} & \textbf{0.1213} & \textbf{0.0757} & 0.1912 & 0.0895 & 0.0589 & \underline{0.1480} & \underline{0.0689} & \underline{0.0452} \\
\bottomrule
\end{tabular}%
}
\end{table}

\section{Ablation Study}

To isolate the contribution of each architectural component, Table~\ref{tab:ablation} compares three configurations on ML-1M: (1) pure BERT4Rec with no auxiliary encoder, (2) BERT-Hybrid-Fixed with a globally shared learnable \(\alpha\), and (3) BERT-Hybrid-Discrete with per-user length-adaptive \(\alpha(u)\) \parencite{ouledmeriem2025hybrid}.

\begin{table}[h]
\centering
\caption{Ablation study on ML-1M (HR@10 / NDCG@10 / MRR@10).}
\label{tab:ablation}
\renewcommand{\arraystretch}{1.3}
\small
\begin{tabular}{lccc}
\toprule
\textbf{Configuration} & \textbf{HR@10} & \textbf{NDCG@10} & \textbf{MRR@10} \\
\midrule
BERT4Rec (no aux. encoder)    & 0.1374 & 0.0652 & 0.0437 \\
+ GNN, fixed \(\alpha\)       & \textbf{0.1485} & \textbf{0.0727} & \textbf{0.0499} \\
+ GNN, adaptive \(\alpha(u)\) & 0.1427 & 0.0682 & 0.0459 \\
\bottomrule
\end{tabular}
\end{table}

Adding the GNN with a fixed fusion weight improves HR@10 by \(+8.1\%\) over BERT4Rec, confirming that global item co-occurrence structure is a genuinely complementary signal \parencite{ouledmeriem2025hybrid}. The length-adaptive variant achieves intermediate overall metrics but delivers superior targeted gains for short-history users, demonstrating that the benefit of adaptive fusion is most pronounced at the user-group level rather than in aggregate statistics.

Table~\ref{tab:best_per_metric} summarizes the best model per metric and user group. BERT-Hybrid-Fixed dominates all overall metrics and medium-history metrics, while TGT-BERT4Rec dominates all short-history metrics across every cutoff.

\begin{table}[h]
\centering
\caption{Best model per metric and user group on ML-1M.}
\label{tab:best_per_metric}
\renewcommand{\arraystretch}{1.3}
\small
\begin{tabular}{llll}
\toprule
\textbf{Group} & \textbf{Metric} & \textbf{Best Model} & \textbf{Value} \\
\midrule
Overall & HR@10   & BERT-Hybrid-Fixed & 0.1485 \\
Overall & NDCG@10 & BERT-Hybrid-Fixed & 0.0727 \\
Overall & MRR@10  & BERT-Hybrid-Fixed & 0.0499 \\
\midrule
Short   & HR@10   & TGT-BERT4Rec      & 0.2778 \\
Short   & NDCG@10 & TGT-BERT4Rec      & 0.1213 \\
Short   & MRR@10  & TGT-BERT4Rec      & 0.0757 \\
\midrule
Medium  & HR@10   & BERT-Hybrid-Fixed & 0.1927 \\
Medium  & NDCG@10 & BERT-Hybrid-Fixed & 0.0963 \\
Medium  & MRR@10  & BERT-Hybrid-Fixed & 0.0672 \\
\bottomrule
\end{tabular}
\end{table}

\section{Efficiency and Statistical Significance}

A key strength of the proposed hybrid family is its \textbf{parameter efficiency}. All four variants add only \(+0.1\)M parameters (\(+5.9\%\)) over BERT4Rec, from 1.7M to 1.8M total trainable parameters \parencite{ouledmeriem2025hybrid}. The additional parameters correspond exclusively to the GCN weight matrices and the projection layer, confirming that performance gains are driven by architectural complementarity rather than increased model capacity.

Statistical significance was assessed using an approximate permutation test \parencite{ouledmeriem2025hybrid}. BERT-Hybrid-Fixed, BERT-Hybrid-Discrete, and TGT-BERT4Rec each achieve statistically significant improvements over BERT4Rec on HR@10 (\(p < 0.05\)). TCN-BERT4Rec shows a smaller, non-significant gain on ML-1M overall, though it achieves the strongest results on ML-100K, suggesting it benefits more from smaller and denser interaction graphs.

% ============================================================
% Experimental Results on MovieLens-100K
% ============================================================
\section{Experimental Results on MovieLens-100K}

This section evaluates the proposed hybrid architectures on the MovieLens-100K benchmark using the same experimental protocol adopted for the larger-scale dataset. The comparison includes the same baseline models: BERT4Rec, BERT-Hybrid-Discrete, BERT-Hybrid-Fixed, Caser, GRU4Rec, LightGCN, SASRec, TCN-BERT4Rec, and TGT-BERT4Rec. The reported results cover global test metrics, sequence-length-aware analysis, and training dynamics.

\subsection{Overall Test Performance}

The global test results on MovieLens-100K indicate that TCN-BERT4Rec delivers the most consistent ranking performance across the majority of evaluation metrics. As shown in Table~\ref{tab:ml100k_test}, TCN-BERT4Rec achieves the best values for HR@5, NDCG@5, HR@10, NDCG@10, HR@20, and NDCG@20, with scores of 0.0682, 0.0386, 0.1194, 0.0551, 0.2111, and 0.0780, respectively. These results confirm that incorporating temporal convolution into the BERT4Rec backbone improves the model’s ability to capture local sequential dependencies while preserving strong contextual modeling capacity.

\begin{table}[h]
\centering
\caption{Overall test metrics on MovieLens-100K}
\label{tab:ml100k_test}
\renewcommand{\arraystretch}{1.3}
\adjustbox{max width=\textwidth}{%
\begin{tabular}{lrrrrrrrrr}
\toprule
\textbf{Model} &
\textbf{HR@5} & \textbf{NDCG@5} & \textbf{MRR@5} &
\textbf{HR@10} & \textbf{NDCG@10} & \textbf{MRR@10} &
\textbf{HR@20} & \textbf{NDCG@20} & \textbf{MRR@20} \\
\midrule
BERT4Rec
  & 0.0544 & 0.0311 & 0.0238
  & 0.0949 & 0.0449 & 0.0294
  & 0.1844 & 0.0671 & 0.0357 \\
BERT-Hybrid-Discrete
  & 0.0578 & 0.0327 & 0.0249
  & 0.1005 & 0.0469 & 0.0305
  & 0.1923 & 0.0698 & 0.0370 \\
BERT-Hybrid-Fixed
  & 0.0595 & 0.0341 & 0.0261
  & 0.1039 & 0.0488 & 0.0319
  & 0.1987 & 0.0725 & 0.0386 \\
Caser
  & 0.0624 & 0.0360 & 0.0279
  & 0.1078 & 0.0504 & 0.0330
  & 0.1994 & 0.0732 & 0.0394 \\
GRU4Rec
  & 0.0341 & 0.0188 & 0.0140
  & 0.0655 & 0.0291 & 0.0181
  & 0.1325 & 0.0458 & 0.0227 \\
LightGCN
  & 0.0227 & 0.0122 & 0.0092
  & 0.0447 & 0.0185 & 0.0114
  & 0.0916 & 0.0293 & 0.0141 \\
SASRec
  & 0.0535 & 0.0305 & 0.0235
  & 0.0959 & 0.0445 & 0.0292
  & 0.1887 & 0.0667 & 0.0355 \\
TCN-BERT4Rec
  & \textbf{0.0682} & \textbf{0.0386} & \textbf{0.0297}
  & \textbf{0.1194} & \textbf{0.0551} & \textbf{0.0362}
  & \textbf{0.2111} & \textbf{0.0780} & \textbf{0.0421} \\
TGT-BERT4Rec
  & \underline{0.0649} & \underline{0.0369} & \underline{0.0283}
  & \underline{0.1140} & \underline{0.0524} & \underline{0.0345}
  & \underline{0.2060} & \underline{0.0748} & \underline{0.0404} \\
\bottomrule
\end{tabular}%
}
\end{table}

On ML-100K, TCN-BERT4Rec achieves the best overall results across all cutoffs, outperforming BERT4Rec by \(+25.8\%\) on HR@10 (0.1194 vs.\ 0.0949). This contrasts with ML-1M, where BERT-Hybrid-Fixed dominates, and suggests that dilated temporal convolutions are particularly effective on smaller and denser interaction graphs. TGT-BERT4Rec is the consistent runner-up, confirming that temporal graph signals provide a complementary advantage over both the GCN and TCN-based auxiliary encoders.

% ============================================================
% CHAPTER 4 — LLM ITEM DESCRIPTION STRATEGIES
% ============================================================
\chapter{LLM Item Description Strategies for Sequential Recommendation}
\addcontentsline{toc}{chapter}{Chapter 4 — LLM Item Description Strategies}
\label{ch:promptcraft}

% ============================================================
\section{Introduction and Motivation}
% ============================================================

The LLM2Sequential framework established by \textcite{boz2025improving}
demonstrated that replacing randomly initialized item embeddings with
dense representations generated by a pretrained Large Language Model
yields consistent and substantial gains in sequential recommendation
quality. The central mechanism is straightforward: each item is
described by a short textual string, that string is passed through
a frozen LLM encoder, and the resulting dense vector initializes
--- and optionally fine-tunes --- the item embedding matrix of a
downstream sequential model such as SASRec or BERT4Rec.

Despite the practical impact of this finding, the original study
leaves a critical design question unanswered: \textbf{how should
an item be described to the LLM?} In the reference experiments,
items are represented primarily by their bare title, with one dataset
supplemented by additional metadata tags. No principled investigation
across multiple formulation strategies is provided, and the authors
themselves acknowledge in passing that enriching the description
with tags improves results on one dataset --- a result whose
implications are never systematically pursued.

This gap matters because the textual description is the sole
information channel between the item's content and the LLM encoder.
The LLM cannot observe interaction patterns, purchase histories, or
co-occurrence statistics; it can only encode what the description
provides. A title-only description encodes little more than the item's
name, while richer formulations --- metadata-enriched attribute lists,
prose narratives built from the product description field,
user-centric summaries framing the item through community preferences,
or comparative profiles anchoring the item to related products ---
may provide the LLM with substantially more discriminative and
contextually relevant information. Each formulation activates a
different subset of the LLM's pretrained knowledge, and the resulting
embeddings may differ significantly in their discriminative capacity
for the recommendation task.

This chapter presents \textbf{Contribution 2} of this memoir: a
systematic investigation of seven item description strategies for
LLM-based sequential recommendation, evaluated on the Amazon Beauty
dataset using BGE-M3 as the embedding model and SASRec as the
downstream sequential model. The chapter is organized as follows.
Section~\ref{sec:context} reviews the LLM2Sequential framework and
its limitations. Section~\ref{sec:pipeline} describes the complete
experimental pipeline implemented in this work. Section~\ref{sec:strategies}
introduces and formalizes the seven description strategies.
Section~\ref{sec:embedding_quality} analyzes the geometric properties
of the resulting embedding spaces. Section~\ref{sec:promptcraft_results}
presents and discusses the experimental results. Section~\ref{sec:promptcraft_summary}
summarizes the key findings.


% ============================================================
\section{Context and Positioning: The LLM2Sequential Framework}
\label{sec:context}
% ============================================================

\subsection{The LLM2Sequential Approach}

\textcite{boz2025improving} systematically evaluated three
orthogonal integration strategies for incorporating LLMs into
sequential recommendation pipelines. Among these, the
\textbf{LLM2Sequential} approach proved to be the most impactful
and the most practically viable. Its core idea is embedding
substitution: the item embedding matrix $\mathbf{E} \in
\mathbb{R}^{n \times d}$ of an existing sequential model
--- which is otherwise randomly initialized and learned
purely from interaction data --- is instead initialized with
dense semantic vectors extracted from a pretrained LLM.

Formally, for each item $i \in \mathcal{I}$, a textual description
$\delta(i)$ is constructed from the item's available metadata.
This description is fed as input to a pretrained LLM encoder
$\phi_{\text{LLM}} : \mathcal{V}^* \rightarrow \mathbb{R}^{d_{\text{LLM}}}$,
where $\mathcal{V}^*$ denotes the space of token sequences and
$d_{\text{LLM}}$ is the LLM's output embedding dimensionality.
The resulting dense vector $\mathbf{z}_i = \phi_{\text{LLM}}(\delta(i))$
is then projected to the model's target dimensionality $d$ via PCA,
yielding the initialization vector $\mathbf{e}_i^{(0)}
\in \mathbb{R}^d$. The downstream sequential model is subsequently
trained normally, either keeping these vectors frozen (feature
extraction mode) or fine-tuning them (warm-start mode).

\textcite{boz2025improving} report that LLM2Sequential initialization
improves NDCG@10 by 15--20\% on their benchmarks compared to
randomly initialized BERT4Rec, providing compelling evidence that
LLM-generated semantic embeddings are highly complementary to the
temporal modeling capabilities of self-attention architectures.

\subsection{The Unresolved Description Strategy Question}

A fundamental limitation of the \textcite{boz2025improving} study is
the absence of any systematic investigation into the description
function $\delta(i)$. In their experiments, $\delta(i)$ defaults to
the item's title string, with one dataset additionally including
user-generated tags. This single-strategy design raises an important
open question: \textbf{is the formulation of $\delta(i)$ a
performance-critical design choice, or does any reasonable description
yield essentially equivalent embeddings?}

The question is theoretically non-trivial. On one hand, the BGE-M3
encoder is trained on hundreds of millions of text pairs and has
been shown to produce high-quality dense representations across a
wide range of text styles and lengths; it might be argued that the
encoder's robustness renders the exact formulation secondary.
On the other hand, different formulations activate different
aspects of the LLM's pretrained knowledge: a title activates
name-level identity information, a category-enriched description
activates taxonomic knowledge, a comparative profile activates
relational knowledge about product similarity, and a user-centric
narrative activates sentiment and preference-related representations.
These different activation patterns may produce embeddings with
fundamentally different geometric structures in the representation
space, with direct consequences for the discriminative capacity of
the resulting item embeddings.

This unresolved question directly motivates the experimental
design of Contribution 2.


% ============================================================
\section{The Complete LLM2Sequential Pipeline}
\label{sec:pipeline}
% ============================================================

\subsection{Dataset: Amazon Beauty}

All experiments in this chapter are conducted on the
\textbf{Amazon Beauty} dataset, a widely used benchmark for
sequential recommendation derived from Amazon product reviews
\parencite{mcauley2015image}. The dataset covers the Beauty
product category and provides user-item interaction sequences
(implicit feedback from purchases and reviews), item titles,
brand names, product categories, descriptions, pricing information,
and co-purchase relationships.

Following the standard preprocessing protocol of
\textcite{kang2018self}, we apply a 5-core filter (retaining only
users and items with at least 5 interactions), sort interactions
chronologically per user, and adopt the \textbf{leave-one-out}
evaluation protocol: the last item of each user's sequence is
held out as the test target, the penultimate item as the validation
target, and all preceding items form the training sequence.

\subsection{Embedding Model: BGE-M3}

We use \textbf{BGE-M3} (BAAI/bge-m3) as the LLM encoder
$\phi_{\text{LLM}}$. BGE-M3 is a multilingual dense retrieval
model developed by the Beijing Academy of Artificial Intelligence,
capable of producing dense representations, sparse lexical
representations, and multi-vector ColBERT-style representations
from the same encoder \parencite{chen2024bgem3}. We extract only
the dense vector output, which has dimensionality $d_{\text{LLM}}
= 1024$.

BGE-M3 is selected for three reasons. First, it represents the
state of the art in general-purpose dense text embedding at the
time of this study, achieving top performance on the MTEB
benchmark. Second, its architecture as a bidirectional encoder
--- in contrast to autoregressive LLMs --- is particularly suited
for producing holistic item-level representations rather than
next-token predictions. Third, its multilingual training corpus
provides robustness to the heterogeneous and sometimes noisy
product text found in the Amazon dataset.

All items are encoded in a single offline pass using
\texttt{fp16} precision on a GPU, exploiting BGE-M3's
\texttt{FlagEmbedding} interface with a maximum input length
of 128 tokens and a batch size of 512. The encoding of the full
item catalog takes approximately 52--56 seconds per strategy,
confirming that the approach is computationally viable.

\subsection{PCA Dimensionality Reduction}

The raw 1024-dimensional BGE-M3 embeddings are reduced to
$d = 256$ dimensions using \textbf{Principal Component Analysis
(PCA)}. A separate PCA transformation is fitted independently
for each description strategy, ensuring that each strategy's
embedding space is optimally compressed along its own principal
directions. This design choice avoids conflating the geometric
properties of different strategies under a shared projection.

PCA is applied \textit{before} feeding the embeddings into
SASRec, replicating the preprocessing protocol of the reference
repository on which this work builds. The number of components
$d = 256$ balances representational capacity with training
efficiency, and is consistent with the downstream model's
hidden dimension.

\subsection{Downstream Model: LLM2SASRec}

The downstream sequential model is \textbf{SASRec}
\parencite{kang2018self} initialized with LLM-derived item
embeddings --- a configuration we term \textbf{LLM2SASRec}.
SASRec is selected as the downstream model for three reasons:
its architectural simplicity allows the experimental results
to be attributed cleanly to the quality of the item embeddings
rather than to the sequential modeling architecture; it is one
of the most widely adopted baselines in sequential recommendation
research; and its unidirectional causal attention mechanism
makes it well-suited to the leave-one-out next-item prediction
evaluation setting.

SASRec is configured uniformly across all strategies with
the following hyperparameters: maximum sequence length $N = 50$,
$L = 1$ Transformer block, $h = 2$ attention heads, embedding
and hidden dimension $d = 256$, dropout rate $p = 0.5$,
learning rate $\eta = 0.001$, weight decay $\lambda = 10^{-4}$,
batch size 256, and early stopping with patience 10 epochs.
The item embedding matrix is initialized from the
strategy-specific PCA-reduced embeddings and fine-tuned during
training.

A \textbf{Vanilla SASRec} baseline --- identical architecture
with randomly initialized (rather than LLM-initialized)
item embeddings --- is trained and evaluated under the same
conditions to quantify the absolute gain attributable to
LLM embedding initialization, independently of description
strategy.

\subsection{Evaluation Protocol}

All models are evaluated under the standard \textbf{leave-one-out}
protocol using the full item set as the candidate pool (no negative
sampling at test time). We report three complementary ranking
metrics: \textbf{NDCG@10}, \textbf{HR@10}, and \textbf{MRR},
following the definitions introduced in Section~1.6 of this memoir.
NDCG@10 is the primary evaluation criterion, as it simultaneously
captures whether the ground-truth item is retrieved and how highly
it is ranked. HR@10 measures coverage at rank 10, and MRR
emphasizes top-1 precision.

The Vanilla SASRec baseline achieves NDCG@10 $= 0.0275$,
HR@10 $= 0.0472$, and MRR $= 0.0229$. All LLM-initialized
strategies are evaluated relative to this baseline to quantify
the gain from semantic initialization, and relative to the
Title-Only strategy (T1) to quantify the incremental gain
from richer description formulations.


% ============================================================
\section{Item Description Strategies}
\label{sec:strategies}
% ============================================================

This section formalizes the seven item description strategies
evaluated in this work. For each item $i$ with metadata
$\mathcal{M}(i)$ containing fields $\{\text{title},
\text{brand}, \text{category}, \text{description},
\text{price}, \text{also\_bought}\}$, we define a description
function $\delta_k(i)$ for $k = 1, \ldots, 7$ that constructs
the textual string passed to the BGE-M3 encoder.

Let us define the following helper extractions from $\mathcal{M}(i)$:
\begin{align*}
  t(i) &= \text{item title (required field)}, \\
  b(i) &= \text{brand name, or ``unknown'' if absent}, \\
  c(i) &= \text{top-3 subcategories (excluding root category)}, \\
  p(i) &= \text{price, or ``unknown'' if absent}, \\
  \text{desc}(i) &= \text{product description text (first 200 chars)}, \\
  R(i) &= \text{top-3 titles from the ``also\_bought'' co-purchase list}.
\end{align*}

\subsection{T1 --- Title Only (Baseline)}

The simplest possible description strategy uses the item title
as the sole input to the LLM encoder:
\begin{equation}
  \delta_1(i) = t(i).
\end{equation}
This strategy directly replicates the default approach of
\textcite{boz2025improving} and serves as the reference
baseline against which all richer formulations are compared.
The title typically contains 3--10 tokens for beauty products
(e.g., \textit{``L'Or\'eal Paris Elvive Total Repair 5 Shampoo''}),
providing the LLM with identity-level information but none of
the contextual signals that characterize the item's function,
audience, or positioning relative to related products.

\subsection{T2 --- Structured Attributes}

The structured attributes strategy assembles the most factually
informative metadata fields into a pipe-separated attribute list:
\begin{equation}
  \delta_2(i) = t(i)
    \;\|\; \bigl[ \texttt{Brand: } b(i) \bigr]
    \;\|\; \bigl[ \texttt{Category: } c(i) \bigr]
    \;\|\; \bigl[ \texttt{Price: } p(i) \bigr],
\end{equation}
where $[\cdot]$ denotes optional inclusion (the field is
appended only if it is non-missing), and $\|$ denotes the
pipe character separator. An example description takes
the form: \textit{``Moroccan Oil Treatment | Brand: OGX
| Category: Hair Care, Serums, Treatments | Price: \$12.99''}.

This strategy provides the LLM with structured taxonomic
and commercial context. The pipe separator is chosen to
signal to the encoder that the fields are independent
attribute-value pairs rather than a flowing narrative,
leveraging the LLM's exposure to structured product listings
during pretraining.

\subsection{T3 --- Rich Prose}

The rich prose strategy constructs a natural language sentence
that integrates the title, brand, category, and product
description into a coherent narrative:
\begin{equation}
  \delta_3(i) =
    t(i)
    \;\text{by}\; b(i) \mathbf{,}
    \;\text{a}\; c(i) \;\text{product.}\;
    \text{desc}(i).
\end{equation}
An example: \textit{``OGX Moroccan Argan Oil Shampoo by OGX,
a Hair Care, Moisturizing Shampoos product. Infused with argan
oil to strengthen and smooth hair\ldots''}.

This strategy is the most information-dense of the seven,
incorporating the product description field which may contain
detailed functional specifications. The hypothesis is that
natural language integration reduces the structural ambiguity
inherent in pipe-separated lists and more closely resembles
the document-style text on which LLMs are pretrained.

\subsection{T4 --- User-Centric}

The user-centric strategy reformulates the item description
from the perspective of the typical consumer, framing the
item through the lens of user preference:
\begin{equation}
  \delta_4(i) = \texttt{``Users who like } t(i)
    \;\texttt{enjoy: } c(i)
    \;\texttt{products from } b(i)\texttt{''}.
\end{equation}
An example: \textit{``Users who like Moroccan Oil Treatment
enjoy: Hair Care, Serums, Treatments products from OGX''}.

This strategy attempts to encode the item's audience profile
directly into the description, activating the LLM's knowledge
of preference patterns and user communities rather than its
knowledge of product properties. The hypothesis is that
preference-aligned language may produce embeddings that are
more closely aligned with the collaborative filtering signals
that sequential models exploit.

\subsection{T5 --- Comparative}

The comparative strategy contextualizes the item within its
neighborhood in the product graph by referencing co-purchased
items:
\begin{equation}
  \delta_5(i) = t(i)
    \;\texttt{is similar to: } R(i)
    \texttt{. Appeals to fans of: } c(i).
\end{equation}
where $R(i)$ is the concatenation of up to three related
product titles derived from the co-purchase graph.
An example: \textit{``Moroccan Oil Treatment is similar to:
Argan Oil Hair Mask, Coconut Oil Conditioner, Keratin Serum.
Appeals to fans of: Hair Care, Serums, Treatments''}.

This strategy encodes relational product knowledge that is
directly aligned with the collaborative filtering objective:
if two items are frequently co-purchased, their embeddings
should be close in the representation space. By explicitly
naming co-purchased products in the description, the strategy
asks the LLM to activate relational semantic associations
rather than purely categorical ones.

\subsection{T6 --- Hybrid}

The hybrid strategy combines the title with the two most
informative structured fields (category and brand) in a
compact pipe-separated format, explicitly omitting the
product description and comparative fields:
\begin{equation}
  \delta_6(i) = t(i)
    \;\|\; \texttt{Category: } c(i)
    \;\|\; \texttt{Brand: } b(i).
\end{equation}
An example: \textit{``Moroccan Oil Treatment | Category:
Hair Care, Serums | Brand: OGX''}.

This strategy is designed as a concise hybrid that provides
categorical and brand context without the verbosity of T3
or the relational dependency on co-purchase data of T5.
It was labeled as the ``best'' strategy in preliminary runs,
motivating its inclusion as a distinct category.

\subsection{T7 --- Structured Comparative}

The structured comparative strategy combines the structured
attribute format of T2 with the co-purchase similarity
signal of T5:
\begin{equation}
  \delta_7(i) = t(i)
    \;\|\; \texttt{Category: } c(i)
    \;\|\; \texttt{Brand: } b(i)
    \;\|\; \texttt{Similar to: } R_2(i),
\end{equation}
where $R_2(i)$ uses up to two related product titles to limit
description length.

This strategy tests whether the combination of structured
metadata and comparative relational context is additive,
or whether the two signals provide redundant information
to the encoder. An example: \textit{``Moroccan Oil Treatment
| Category: Hair Care, Serums | Brand: OGX | Similar to:
Argan Oil Hair Mask, Coconut Oil Conditioner''}.

\subsection{Summary of Strategy Design Space}

Table~\ref{tab:strategies_summary} summarizes the design
dimensions of the seven strategies.

\begin{table}[ht]
\centering
\caption{Design dimensions of the seven description strategies}
\label{tab:strategies_summary}
\renewcommand{\arraystretch}{1.3}
\footnotesize
\begin{tabularx}{\textwidth}{@{} l c c c c c @{}}
\toprule
\textbf{Strategy} & \textbf{Title} & \textbf{Brand} &
\textbf{Category} & \textbf{Description} & \textbf{Co-purchase} \\
\midrule
T1: Title Only          & \checkmark & $\times$ & $\times$ & $\times$ & $\times$ \\
T2: Structured Attrs    & \checkmark & \checkmark & \checkmark & $\times$ & $\times$ \\
T3: Rich Prose          & \checkmark & \checkmark & \checkmark & \checkmark & $\times$ \\
T4: User-Centric        & \checkmark & \checkmark & \checkmark & $\times$ & $\times$ \\
T5: Comparative         & \checkmark & $\times$ & \checkmark & $\times$ & \checkmark \\
T6: Hybrid              & \checkmark & \checkmark & \checkmark & $\times$ & $\times$ \\
T7: Struct+Compare      & \checkmark & \checkmark & \checkmark & $\times$ & \checkmark \\
\bottomrule
\end{tabularx}
\end{table}


% ============================================================
\section{Embedding Space Analysis}
\label{sec:embedding_quality}
% ============================================================

Before examining downstream recommendation quality, we
analyze two geometric properties of the raw 1024-dimensional
BGE-M3 embeddings produced by each strategy: \textbf{isotropy}
and \textbf{average pairwise cosine distance}. These
intrinsic measures characterize the structure of the
embedding space independently of any downstream model.

\subsection{Isotropy}

Embedding isotropy measures how uniformly distributed the
embedding vectors are across all directions of the
representation space \parencite{mu2018allbutthetop}. A
perfectly isotropic embedding space uses every dimension
equally, while an anisotropic space concentrates variance
along a small number of principal directions, wasting
representational capacity. Formally, isotropy is measured
as the ratio of the smallest to the largest singular value
of the centered embedding matrix:
\begin{equation}
  \text{Isotropy} = \frac{\sigma_{\min}}{\sigma_{\max}},
\end{equation}
where $\sigma_{\min}$ and $\sigma_{\max}$ are respectively
the smallest and largest singular values of
$\mathbf{Z} - \bar{\mathbf{z}}\mathbf{1}^\top$, with
$\mathbf{Z} \in \mathbb{R}^{n \times d_{\text{LLM}}}$ the
embedding matrix and $\bar{\mathbf{z}}$ the mean embedding.

In our experiments, all seven strategies yield an isotropy
score of approximately 0.0, indicating that the BGE-M3
embeddings are highly anisotropic for the Beauty item catalog.
This result is consistent with the well-documented
\textbf{representation degeneration} phenomenon in contextual
embedding models \parencite{li2020sentence, mu2018allbutthetop},
where pretrained LLM embeddings concentrate in a narrow
cone of the representation space. The dominant singular value
is orders of magnitude larger than the smallest, reflecting
that most semantic variation is captured by a small number
of principal directions --- which is precisely why PCA
reduction to 256 dimensions retains a large fraction of
the total variance.

\subsection{Average Pairwise Cosine Distance}

Average pairwise cosine distance measures the typical
dissimilarity between item embeddings, providing an estimate
of how \textit{discriminative} the embedding space is:
higher values indicate that the strategy produces more
spread-out embeddings where different items are more
easily separated. We estimate this quantity from a random
sample of 2\,000 item pairs:
\begin{equation}
  \overline{d}_{\cos} =
    1 - \frac{1}{\binom{2000}{2}}
    \sum_{i \neq j} \cos(\mathbf{z}_i, \mathbf{z}_j).
\end{equation}

Table~\ref{tab:embedding_quality} reports the average
pairwise cosine distance for each strategy.

\begin{table}[ht]
\centering
\caption{Embedding space geometry: average pairwise cosine distance per strategy (higher = more discriminative)}
\label{tab:embedding_quality}
\renewcommand{\arraystretch}{1.3}
\small
\begin{tabular}{@{} l c @{}}
\toprule
\textbf{Strategy} & \textbf{Avg.\ Pairwise Cosine Dist.} \\
\midrule
T1: Title Only (Baseline) & 0.5801 \\
T2: Structured Attrs      & 0.4868 \\
T3: Rich Prose            & 0.5273 \\
T4: User-Centric          & 0.4229 \\
T5: Comparative           & 0.4814 \\
T6: Hybrid                & 0.5256 \\
T7: Struct+Compare        & 0.5080 \\
\bottomrule
\end{tabular}
\end{table}

A clear and somewhat counterintuitive pattern emerges from
Table~\ref{tab:embedding_quality}: \textbf{T1 (Title Only)
produces the most spread-out embedding space} ($\overline{d}_{\cos}
= 0.5801$), while \textbf{T4 (User-Centric) produces the most
compact space} ($\overline{d}_{\cos} = 0.4229$). The richer
structured and comparative strategies (T2, T5, T7) cluster
in the middle range 0.48--0.51.

This pattern can be explained as follows. A bare item title
is typically the most uniquely identifying text for a product:
two beauty items with different names will produce maximally
distinct title-only descriptions, activating very different
regions of the LLM's embedding space. In contrast, the
user-centric formulation \textit{``Users who like X enjoy:
[category] products from [brand]''} introduces a common
grammatical and semantic template across all items in the same
category-brand combination, pulling their embeddings toward
a shared semantic center. The comparative strategy (T5)
shows a similar compaction effect due to the shared relational
vocabulary (``is similar to'', ``appeals to fans of'').

Crucially, this analysis reveals that \textbf{geometric
discriminability of the embedding space is not a reliable
proxy for downstream recommendation quality}: T1 has the
highest pairwise distance yet ranks only as the reference
baseline in NDCG@10, while T5 --- despite a notably
lower pairwise distance of 0.4814 --- achieves the best
recommendation performance. We interpret this as evidence
that the relevant notion of discriminability for recommendation
is not global spread across all item pairs, but rather
local discriminability between items that are semantically
distinct from the user's perspective --- a distinction
that comparative and relational descriptions may capture
more effectively.


% ============================================================
\section{Experimental Results and Analysis}
\label{sec:promptcraft_results}
% ============================================================

\subsection{Main Results}

Table~\ref{tab:promptcraft_main} reports the full experimental
results for all seven description strategies and the Vanilla
SASRec baseline on the Amazon Beauty dataset.

\begin{table}[ht]
\centering
\caption{PromptCraft-SeqRec results on Amazon Beauty. All LLM-initialized strategies vs.\ Vanilla SASRec and T1 baseline. Best LLM result per metric in \textbf{bold}, second best \underline{underlined}.}
\label{tab:promptcraft_main}
\renewcommand{\arraystretch}{1.3}
\small
\begin{adjustbox}{max width=\textwidth}
\begin{tabular}{@{} l c c c c c c @{}}
\toprule
\textbf{Strategy} &
\textbf{NDCG@10} & \textbf{NDCG@20} &
\textbf{HR@10} & \textbf{HR@20} &
\textbf{MRR} &
\textbf{$\Delta$NDCG@10 vs T1} \\
\midrule
Vanilla SASRec (no LLM)   & 0.0275 & 0.0324 & 0.0472 & 0.0667 & 0.0229 & $-21.5\%$ \\
\midrule
T1: Title Only (Baseline) & 0.0351 & 0.0423 & 0.0642 & 0.0931 & 0.0281 & --- \\
T2: Structured Attrs      & \underline{0.0359} & \underline{0.0438} & 0.0635 & \underline{0.0949} & \textbf{0.0297} & $+2.3\%$ \\
T3: Rich Prose            & 0.0340 & 0.0411 & 0.0618 & 0.0902 & 0.0275 & $-3.0\%$ \\
T4: User-Centric          & 0.0357 & 0.0433 & 0.0631 & 0.0935 & \underline{0.0294} & $+1.7\%$ \\
T5: Comparative           & \textbf{0.0360} & \textbf{0.0438} & \underline{0.0658} & \textbf{0.0967} & 0.0291 & $+\mathbf{2.6\%}$ \\
T6: Hybrid                & 0.0354 & 0.0427 & 0.0635 & 0.0926 & 0.0289 & $+0.8\%$ \\
T7: Struct+Compare        & 0.0359 & 0.0432 & \textbf{0.0671} & \underline{0.0962} & 0.0285 & $+2.4\%$ \\
\bottomrule
\end{tabular}
\end{adjustbox}
\end{table}

\subsection{LLM Initialization vs.\ Vanilla SASRec}

A first and unambiguous finding is that \textbf{all seven
LLM-initialized strategies outperform the Vanilla SASRec
baseline} (NDCG@10 $= 0.0275$) by a substantial margin.
The gain ranges from $+23.6\%$ (T3 Rich Prose,
NDCG@10 $= 0.0340$) to $+30.8\%$ (T5 Comparative,
NDCG@10 $= 0.0360$), confirming the robustness of the
LLM2Sequential initialization strategy across all
description formulations. Even the weakest strategy (T3)
yields a substantial absolute improvement, establishing
that semantic LLM initialization is beneficial regardless
of the description quality.

This result is aligned with and replicates the findings of
\textcite{boz2025improving} on a different dataset (Amazon
Beauty vs.\ their MovieLens benchmarks), providing additional
cross-domain evidence for the generalizability of the
LLM2Sequential approach.

\subsection{Impact of Description Strategy on NDCG@10}

The central question of this chapter is whether the description
formulation affects recommendation quality. The results in
Table~\ref{tab:promptcraft_main} provide a clear affirmative
answer: the gap between the best and worst LLM strategies
(T5 vs.\ T3) amounts to $+5.9\%$ relative NDCG@10, a
non-trivial difference that is larger than the gains typically
reported by architectural innovations in the sequential
recommendation literature.

The ranking of strategies by NDCG@10 is as follows:
\[
  \text{T5} \approx \text{T2} \approx \text{T7} > \text{T4} > \text{T6} > \text{T1} \gg \text{T3}.
\]
The top three strategies (T5, T2, T7) all improve over the
title-only baseline by approximately $+2.3\%$ to $+2.6\%$
NDCG@10. T4 (User-Centric) provides a smaller but consistent
$+1.7\%$ improvement. T6 (Hybrid) is marginally above the
baseline ($+0.8\%$). T3 (Rich Prose) is the only strategy
to degrade performance, falling $3.0\%$ below the title-only
baseline.

\subsection{Metric-Specific Analysis}

The three metrics tell a slightly different story at the
level of the top strategies:

\begin{itemize}
  \item \textbf{NDCG@10} (ranking quality): T5 Comparative
        achieves the highest score (0.0360), closely followed
        by T2 Structured and T7 Struct+Compare (0.0359).
        T5's superiority on this position-sensitive metric
        suggests that relational comparative descriptions
        help the model rank the ground-truth item toward
        the top of the list.
  \item \textbf{HR@10} (coverage): T7 Struct+Compare
        achieves the highest HR@10 (0.0671), followed by
        T5 Comparative (0.0658). The combination of
        structured attributes and co-purchase context
        appears to maximize the probability that the correct
        item appears anywhere in the top-10 list.
  \item \textbf{MRR} (top-1 precision): T2 Structured Attrs
        achieves the highest MRR (0.0297), indicating that
        factual structured metadata is particularly effective
        at pushing the ground-truth item to the very top
        position of the ranking.
\end{itemize}

No single strategy dominates across all three metrics,
suggesting that the optimal formulation depends on the
specific recommendation objective. For applications where
top-1 precision is critical (e.g., single-item recommendations),
T2 structured attributes is recommended; for applications
requiring maximum coverage (e.g., top-$K$ diversified lists),
T7 or T5 are preferable.

\subsection{Why Does T3 (Rich Prose) Underperform?}

The underperformance of T3 is the most striking and
counterintuitive finding of this study. T3 is the only
strategy that incorporates the full product description
text (up to 200 characters) and constructs a grammatically
coherent prose narrative. One might expect this richest
formulation to produce the best embeddings. Instead, it
is the only strategy to fall below the title-only baseline.

We attribute this to two complementary factors. First,
product descriptions in the Amazon Beauty dataset are
frequently \textbf{marketing copy} rather than informative
factual content: they contain superlatives, promotional
language, and generic claims (e.g., \textit{``transforms
your hair''}, \textit{``leaves skin glowing''}) that
provide the LLM with noisy, low-discriminative signals.
The BGE-M3 encoder, trained primarily on informative text
pairs, may actually perform worse when presented with
ambiguous or hyperbolic input. Second, the 128-token
maximum length of the BGE-M3 encoding window means that
a significant portion of longer prose descriptions
is \textbf{truncated}, potentially breaking sentence
structure in a way that degrades embedding quality.

\subsection{Why Does T5 (Comparative) Perform Best?}

T5 achieves the best NDCG@10 by incorporating the
co-purchase neighborhood of each item (``also bought''
lists). We argue that this strategy is uniquely aligned
with the sequential recommendation objective: \textbf{items
that users frequently co-purchase are exactly the items
that should be nearby in the recommendation model's
embedding space}. By explicitly naming co-purchased
products in the description, T5 asks the LLM to produce
an embedding that is semantically adjacent to those products
--- directly encoding the collaborative filtering signal
into the embedding initialization before any training
occurs. This pre-alignment may provide the downstream
SASRec model with a warm-start embedding space that is
already partially optimized for the recommendation task.


% ============================================================
\section{Summary and Key Findings}
\label{sec:promptcraft_summary}
% ============================================================

This chapter presented a systematic investigation of seven
item description strategies for LLM-based sequential
recommendation, evaluated on the Amazon Beauty dataset using
BGE-M3 embeddings and SASRec as the downstream model.
The main findings are as follows.

\medskip
\noindent\textbf{Finding 1 --- LLM initialization is robustly beneficial.}
All seven description strategies outperform the Vanilla
SASRec baseline by $+23\%$ to $+31\%$ relative NDCG@10,
confirming the generalizability of the LLM2Sequential
approach across description formulations and datasets.

\medskip
\noindent\textbf{Finding 2 --- Description strategy is a
performance-critical design choice.}
The gap between the best (T5: $+2.6\%$) and worst (T3: $-3.0\%$)
LLM strategies relative to the title-only baseline
amounts to $5.9\%$ relative NDCG@10. This gap exceeds the
gains typically reported by architectural modifications in
the recommendation literature, establishing description
strategy as a first-class design dimension in LLM-augmented
recommendation systems.

\medskip
\noindent\textbf{Finding 3 --- Comparative descriptions best
align with the recommendation objective.}
T5 (Comparative) achieves the best NDCG@10 and T7
(Struct+Compare) achieves the best HR@10, both leveraging
co-purchase relationships. Incorporating relational
product knowledge into the description directly pre-aligns
the embedding space with the collaborative filtering
signal exploited by the sequential model.

\medskip
\noindent\textbf{Finding 4 --- Geometric discriminability
is not a reliable predictor of downstream quality.}
T1 (Title Only) produces the most geometrically spread-out
embeddings ($\overline{d}_{\cos} = 0.5801$) yet ranks only
as the reference baseline. Global pairwise distance is a
poor surrogate for recommendation-relevant discriminability.

\medskip
\noindent\textbf{Finding 5 --- Prose descriptions with noisy
marketing text hurt performance.}
T3 (Rich Prose) is the only strategy to degrade below the
title-only baseline, underscoring the risk of incorporating
unfiltered product description fields that contain
marketing language rather than factually discriminative content.

\begin{table}[ht]
\centering
\caption{Summary of strategy recommendations by objective}
\label{tab:strategy_recommendations}
\renewcommand{\arraystretch}{1.3}
\small
\begin{tabular}{@{} l l @{}}
\toprule
\textbf{Objective} & \textbf{Recommended Strategy} \\
\midrule
Maximize NDCG@10 (ranking quality)  & T5: Comparative \\
Maximize HR@10 (coverage)            & T7: Struct+Compare \\
Maximize MRR (top-1 precision)       & T2: Structured Attrs \\
Minimal metadata (title+category only) & T6: Hybrid \\
Avoid: noisy description fields      & T3: Rich Prose \\
\bottomrule
\end{tabular}
\end{table}

% ============================================================
% GENERAL CONCLUSION
% ============================================================
\chapter*{General Conclusion}
\addcontentsline{toc}{chapter}{General Conclusion}

\section*{Summary of Contributions}

This memoir presented two complementary contributions to the field of sequential recommendation.

\medskip
\noindent\textbf{Contribution 1 --- Length-Adaptive Hybrid GNN-Transformer.}
We proposed a family of four hybrid architectures combining a BERT4Rec Transformer backbone with a complementary auxiliary encoder, unified by a length-adaptive fusion mechanism. The key insight is that the relative utility of global graph signals versus local sequential signals varies systematically with user interaction history length: sparse users benefit most from GNN-based global collaborative signals, while active users benefit from the full expressiveness of the Transformer. The four variants differ in their auxiliary encoder (GCN co-occurrence graph, dilated convolutions, or temporal graph) and fusion strategy (global learnable $\alpha$, per-user discrete $\alpha(u)$, or gated sigmoid). Evaluated on MovieLens-100K and MovieLens-1M, BERT-Hybrid-Fixed achieves the strongest aggregate performance on ML-1M (\(+8.1\%\) HR@10, \(+11.5\%\) NDCG@10, \(+14.2\%\) MRR@10 over BERT4Rec), while TGT-BERT4Rec achieves the highest HR@10 for short-history users. All four variants add only $+0.1$M parameters ($+5.9\%$) over BERT4Rec, confirming that the gains stem from architectural complementarity rather than increased capacity.

\medskip
\noindent\textbf{Contribution 2 --- Systematic Study of LLM Item Description Strategies.}
We conducted the first systematic investigation of how the formulation of item textual descriptions affects the quality of LLM-generated embeddings in the context of sequential recommendation, following the LLM2Sequential framework of \textcite{boz2025improving}. Seven description strategies were designed and evaluated on the Amazon Beauty dataset using BGE-M3 embeddings and SASRec as the downstream model. The comparative strategy (T5) achieved the best NDCG@10 ($+2.6\%$ over the title-only baseline), while Rich Prose (T3) was the only strategy to hurt performance ($-3.0\%$). All LLM-initialized strategies outperformed Vanilla SASRec by $+23\%$ to $+31\%$ relative NDCG@10, and description strategy was shown to be a performance-critical design choice with a 5.9\% gap between best and worst formulations.

\section*{Limitations}

Several methodological limitations should be acknowledged. First, the asymmetric maximum sequence length (200 for hybrid models vs.\ 50 for pure sequential baselines) introduces a confound in the comparison; future work should control for uniform truncation. Second, the discrete bins in HybridDiscrete introduce hard boundary discontinuities and receive no gradient signal; a fully differentiable continuous fusion function would be a natural extension. Third, the evaluation is conducted exclusively on MovieLens benchmarks; generalization to other domains and interaction types remains to be validated.

\section*{Future Directions}

Several promising research directions emerge from this work. First, replacing the discrete length bins with a continuous MLP-based $\alpha(u)$ function trained end-to-end would eliminate boundary effects and allow the model to learn a smooth, data-driven fusion schedule. Second, extending the systematic study of LLM description strategies to more recent instruction-tuned models (GPT-4, LLaMA-3) and to non-movie domains (e-commerce, music) would broaden the practical impact. Third, combining both contributions --- length-adaptive hybrid fusion with LLM-enriched item embeddings --- represents a natural next step toward a unified architecture that leverages structural, temporal, and semantic item signals simultaneously.

% ============================================================
% BIBLIOGRAPHY
% ============================================================
\newpage
\bibliography{references}

\end{document}